\chapter{BV-BRST formalism and Gaiotto's perspective}
\label{ch:bv-brst}
\index{Batalin--Vilkovisky|see{BV algebra}}

The physicist's BRST cohomology and the algebraist's bar resolution
are not analogous: they are the same computation. The BV formalism
is a machine for computing derived functors; the bar complex is a
machine for computing derived functors; at genus~$0$, they are the
same machine applied to the same input. Costello's formulation of
perturbative QFT as the theory of derived moduli problems makes the
identification inevitable: the BRST operator computes the derived
space of gauge-equivalence classes, which is what the bar differential
computes on the algebraic side. What requires proof is not the
coincidence but the precise scope of the agreement.

The BV Laplacian and the sewing operator agree at genus~$0$. The
quantum master equation $\hbar \Delta S + \tfrac12\{S,S\} = 0$ of
the BV formalism coincides on $\mathbb{P}^1$ with the bar
differential $d_{\mathrm{bar}}$ of the chiral algebra: the
BRST operator $Q_{\mathrm{BRST}}$ is the bar differential, the
classical master equation is $d_{\mathrm{bar}}^2 = 0$, and anomaly
cancellation ($c = 26$) is the condition $\kappa_{\mathrm{tot}} = 0$
on the total Virasoro central charge of matter plus ghosts. Do
they agree at all genera?

A natural obstruction sits in the way. At higher genus, the BV
Laplacian receives contributions from handle-gluing amplitudes that
involve the full moduli space $\overline{\mathcal{M}}_{g,n}$, while
the bar differential collects OPE residues along FM boundary strata.
What is proved in this chapter is the genus-$0$ BV/bar comparison
and the all-genera scalar Heisenberg identity, together with the
all-genera coderived comparison of
Theorem~\ref{thm:bv-bar-coderived}. Beyond genus~$0$, the
obstruction analysis isolates why classes~$\mathsf{G}$,
$\mathsf{L}$, and~$\mathsf{C}$ admit stronger chain-level statements
than class~$\mathsf{M}$. The higher-genus comparison is resolved in
$D^{\mathrm{co}}$ because the harmonic discrepancy factors through
the curvature and the resulting cone is $m_0$-torsion, hence
coacyclic. The stronger chain-level statement is unconditional for
classes~$\mathsf{G}$ and~$\mathsf{L}$, conditional for
class~$\mathsf{C}$ on harmonic decoupling, and open for
class~$\mathsf{M}$. The Heisenberg case is resolved at the
scalar level at all genera:
$F_g^{\mathrm{BV}}(\cH_\kappa) = F_g^{\mathrm{bar}}(\cH_\kappa)
= \kappa \cdot \lambda_g^{\mathrm{FP}}$
\textup{(}Theorem~\ref{thm:heisenberg-bv-bar-all-genera}\textup{)}.
For general algebras what survives at every genus is the scalar
anomaly package: the universal MC class $\Theta_\cA$ encodes the
one-loop anomaly coefficient $\kappa(\cA)$, and its projection
onto $H^2(\overline{\mathcal{M}}_{g,n})$ gives $F_g = \kappa(\cA) \cdot \lambda_g$
on the uniform-weight lane \textup{(}Theorem~D\textup{)}. The
multi-weight correction $\delta F_g^{\mathrm{cross}}$ is where the
chain-level identification breaks: the cross-channel term in the
sewing expansion has no counterpart in the BV side until
harmonic decoupling is assumed, and class~$\mathsf{M}$ is precisely
where that assumption fails.

\begin{remark}[BRST anomaly brackets and transferred SC operations]
\label{rem:brst-anomaly-gkw}
\index{Gaiotto--Kulp--Wu!BRST anomaly brackets}
\index{BRST operator!GKW higher operations}
The BRST-anomaly brackets of Gaiotto--Kulp--Wu~\cite{GKW2025}
and the transferred $\mathrm{SC}^{\mathrm{ch,top}}$ operations of
the bar complex are two constructions of the same $A_\infty$
structure: GKW compute higher operations from Wess--Zumino
consistency on a disk chart; the bar complex computes them from
iterated residues on FM boundary strata
(Remark~\ref{rem:iterated-residues-ainfty}).
The MC equation $D \cdot \Theta + \tfrac{1}{2}[\Theta, \Theta] = 0$
(Theorem~\ref{thm:mc2-bar-intrinsic}) is the all-genera master
equation of which GKW's consistency conditions are the genus-$0$
projections.
\end{remark}

\begin{remark}[BV/bar dictionary; \ClaimStatusHeuristic]
\label{rem:bv-bar-bridge}
\index{BV algebra!bar complex bridge|textbf}
Let $\cA$ be a chiral algebra on a smooth curve~$X$. The BV and
bar-cobar formalisms are related by the following dictionary:
\begin{enumerate}[label=\textup{(\roman*)}]
\item the geometric BV complex and the geometric bar complex are
 identified at genus~$0$
 \textup{(}Theorem~\textup{\ref{thm:bv-bar-geometric}}, proved
 elsewhere\textup{)};
\item the classical master equation corresponds to the
 genus-$0$ square-zero condition for the bar differential after
 transporting through that external comparison;
\item the scalar genus-$1$ obstruction is the modular characteristic
 $\kappa(\cA)$, the one-loop anomaly
 coefficient.
\end{enumerate}
The chain-level identifications (antifields versus
cogenerators, the BV Laplacian, the full quantum master equation,
BRST versus bar cohomology) remain heuristic in this manuscript.

Item~\textup{(i)} is Theorem~\ref{thm:bv-bar-geometric}. The scalar
genus-$1$ coefficient in item~\textup{(iii)} is the content of
Theorem~\ref{thm:genus-universality} and, for Virasoro, of
Computation~\ref{comp:virasoro-curvature}. The remaining entries are
the intended BV reading of these proved algebraic statements, not
chain-level equalities.
\end{remark}

\begin{remark}[Heuristic BV reading of the modular MC hierarchy]
\label{rem:modular-qme-bv}
\index{quantum master equation!modular}
\index{modular quantum master equation|textbf}

The proved object is the universal Maurer--Cartan class~$\Theta_\cA$
\textup{(}Theorem~\ref{thm:mc2-bar-intrinsic}\textup{)}.
Write $\Theta_\cA = \Theta_0 + \hbar\,\Theta_1 + \hbar^2\,\Theta_2
+ \cdots$ for the genus expansion. The full modular QME
\begin{equation}\label{eq:modular-qme-bv}
d\,\Theta_\cA + \tfrac{1}{2}[\Theta_\cA, \Theta_\cA]
+ \hbar\,\Delta\,\Theta_\cA = 0
\end{equation}
in the completed cyclic deformation Lie algebra
$\Defcyc^{\mathrm{mod}}(\cA)$ breaks into a hierarchy:
\begin{alignat*}{2}
&\text{genus } 0\colon\quad
 &&d\,\Theta_0 + \tfrac{1}{2}[\Theta_0,\Theta_0] = 0
 \qquad\text{(classical ME $=$ $d_{\mathrm{bar}}^2 = 0$)}, \\
&\text{genus } 1\colon\quad
 &&d\,\Theta_1 + [\Theta_0,\Theta_1] + \Delta\,\Theta_0 = 0
 \qquad\text{(one-loop anomaly $=$ $\kappa(\cA)\cdot\omega_1$)}, \\
&\text{genus } 2\colon\quad
 &&d\,\Theta_2 + [\Theta_0,\Theta_2]
 + \tfrac{1}{2}[\Theta_1,\Theta_1]
 + \Delta\,\Theta_1 = 0
 \qquad\text{(two-loop interference)}.
\end{alignat*}
The displayed hierarchy is a heuristic BV translation of the proved
modular MC hierarchy, not a proved BV identity: the
genus-$0$ line is a classical master equation, the genus-$1$
line a one-loop anomaly equation, and higher lines encode
nonlinear modular interference. At the scalar level, the trace
of~\eqref{eq:modular-qme-bv} recovers the genus tower
$F_g = \kappa(\cA)\lambda_g$
\textup{(}Theorem~D, on the uniform-weight lane\textup{)}. At the
full level, equation~\eqref{eq:modular-qme-bv} is the principal
conjectural target of the modular cumulant transform
(\S\ref{subsec:completion-kinematics-programme} of the concordance).

In the Feynman-diagram language of
Chapter~\ref{ch:feynman}: $\Theta_0$ collects tree
diagrams, $\Theta_1$ collects one-loop diagrams, and $\Theta_g$
collects connected $g$-loop diagrams. In physics language, the modular MC hierarchy is an all-loop Ward identity. The connected
stable-graph exponential
$\Theta_\cA = \sum_{\Gamma\,\text{connected}}
\frac{\hbar^{g(\Gamma)}}{|\operatorname{Aut}(\Gamma)|}\,\Phi_\Gamma$
makes this precise: the vertices are cumulant primitives, internal
edges are factorization propagators, and the symmetry factor is the
standard graph automorphism weight.
\end{remark}

\begin{theorem}[Curved BV quantum master equation at genus $g\ge 1$;
  \ClaimStatusProvedHere]
\label{thm:curved-bv-qme-genus-g}
\index{BV algebra!curved QME}
\index{quantum master equation!curved genus $g$}
Let $\cA$ be a chirally Koszul chiral algebra on a smooth projective
curve of genus $g$, and let $S \in \Defcyc^{\mathrm{mod}}(\cA)$ denote
a solution of the classical master equation on the uniform-weight lane.
The genus-$g$ curved BV quantum master equation in the chiral ambient
reads
\begin{equation}\label{eq:curved-bv-qme}
\tfrac{1}{2}\{S,S\} + \hbar\,\Delta\,S + \hbar^{2}\,\kappa(\cA)\cdot
\omega_g \;=\; 0
\end{equation}
where $\{{-},{-}\}$ is the BV antibracket, $\Delta$ is the BV Laplacian
of \autoref{def:bv-laplacian}, and $\omega_g \in H^{2}(\Mbar_{g,n},\Q)$
is the Arakelov curvature class with $\int_{\Mbar_g}\omega_g = 1$.
\begin{enumerate}[label=\textup{(\roman*)}]
\item The $\hbar^{2}$ term is the genus-$g$ one-loop Arakelov anomaly;
      it is precisely the bar curvature $m_0 = \kappa(\cA)\cdot\omega_g$
      of the bar differential (\autoref{rem:anomaly-curvature-bv},
      Theorem~D).
\item The curved QME \eqref{eq:curved-bv-qme} is solvable iff
      $\kappa(\cA)\cdot\omega_g$ is exact in
      $H^{2}(\Mbar_{g,n},\obs_{\cA})$; this is the modular-bootstrap
      $H^{2}=0$ vanishing of
      \autoref{prop:modular-bootstrap-to-curved-dunn-bridge}.
\item The anomaly-completed class $\Theta_g \in Z^{2}(\cA^{!}_{\partial})$
      of~\autoref{app:anomaly-completed-topological-holography}
      is the closed representative of $\kappa(\cA)\cdot\omega_g$ in the
      Koszul-dual cyclic complex; under the Verdier identification
      $\cA^{!}_{\infty} \simeq \cA^{!}$ on the Koszul locus,
      $\Theta_g$ equals the BV curvature of
      \eqref{eq:curved-bv-qme}.
\end{enumerate}
\end{theorem}

\begin{proof}
\emph{Step~1 (bar $=$ BV curvature).} By
\autoref{thm:bv-bar-geometric} the genus-$0$ classical master equation
$\{S,S\}=0$ is the square-zero bar differential. In the coderived
ambient $D^{\mathrm{co}}$, \autoref{thm:bv-bar-coderived} extends the
bar/BV identification to all genera: the genus-$g$ bar differential
satisfies $d_{\barB}^{\,2} = \kappa(\cA)\cdot\omega_g$
(\autoref{thm:anomaly-koszul} and Theorem~D on the uniform-weight
lane). Under the chain-level isomorphism $C_{\mathrm{BV}}\cong
\barB^{\mathrm{ch}}$ this transfers to
$Q_{\mathrm{BV}}^{2} = \kappa(\cA)\cdot\omega_g$, i.e.\ the curvature
of the BV differential is $m_0 = \kappa(\cA)\cdot\omega_g$.

\emph{Step~2 (curved BV identity).} Expanding
$Q_{\mathrm{BV}} = \{S,{-}\} + \hbar\,\Delta$ on
$C_{\mathrm{BV}}$ and applying $Q_{\mathrm{BV}}^{2}$ to the unit, the
graded Leibniz and $\Delta^{2}=0$ (\autoref{thm:config-space-bv},
which is unconditional at the level of the scalar uniform-weight
projection used here) yield
$Q_{\mathrm{BV}}^{2}(1) = \tfrac{1}{2}\{S,S\} + \hbar\,\Delta S$
plus the $\hbar^{2}$ two-loop tadpole.  Equating with
$m_0 = \kappa(\cA)\cdot\omega_g$ gives
\eqref{eq:curved-bv-qme} on the uniform-weight scalar lane.

\emph{Step~3 (solvability $\Leftrightarrow$ modular bootstrap
$H^{2}=0$).} The curvature class
$\kappa(\cA)\cdot\omega_g$ lies in
$H^{2}(\Mbar_{g,n},\obs_{\cA})$; solving
\eqref{eq:curved-bv-qme} means producing a genus-$g$ effective
action $S_g$ whose curved MC equation absorbs this class. The
modular-bootstrap-to-curved-Dunn bridge
(\autoref{prop:modular-bootstrap-to-curved-dunn-bridge}) supplies a
chain map $\Phi$ whose image of the modular-bootstrap cohomology
$H^{2}_{\mathrm{mod}}=0$ kills the curvature in the curved-Dunn
twisting-cochain complex, providing the required primitive.

\emph{Step~4 (anomaly-completed $\Theta_g$).} In
\autoref{app:anomaly-completed-topological-holography}, $\Theta_g$ is
the central degree-$2$ transgression generator of the anomaly-completed
Koszul-dual algebra $A^{!}_{\Theta}$ satisfying
$d\eta = \Theta_g$ (L1613). Under Verdier duality
$\cA^{!}_{\infty}\simeq \cA^{!}$ on the Koszul locus,
$\Theta_g$ is the image of $\kappa(\cA)\cdot\omega_g$ in
$Z^{2}(\cA^{!}_{\partial})$; the equality is at cohomology in
$D^{\mathrm{co}}$ and strict on the uniform-weight scalar lane.

Steps~1--4 establish \eqref{eq:curved-bv-qme} together with the
Verdier identification $\Theta_g = \kappa(\cA)\cdot\omega_g$ and the
equivalence of solvability with modular-bootstrap
$H^{2}$-vanishing.
\end{proof}

\begin{remark}[Costello--Gwilliam one-loop renormalisation
  $=$ modular bootstrap]
\label{rem:costello-gwilliam-modular-bootstrap}
\index{Costello--Gwilliam!one-loop renormalisation}
The classical Costello--Gwilliam package identifies the genus-$g$
renormalisation obstruction of a perturbative QFT with the
cohomology class of its one-loop anomaly in
$H^{2}(\Mbar_{g,n},\obs)$. Theorem~\ref{thm:curved-bv-qme-genus-g}
identifies that class with $\kappa(\cA)\cdot\omega_g$ for the chiral
boundary $\cA$, and the solvability criterion (ii) with the
modular-bootstrap vanishing
$H^{2}_{\mathrm{mod}}=0$ proved in
\autoref{thm:curved-dunn-H2-vanishing-all-genera}. Thus
Costello--Gwilliam one-loop renormalisability at genus $g\ge 1$
for chirally Koszul $\cA$ on the uniform-weight lane reduces to the
curved bar-cobar inversion of
\autoref{thm:bar-cobar-isomorphism-main}. The anomaly-completed
Yangian curvature $\Theta_g$ of
\autoref{app:anomaly-completed-topological-holography} is the
algebraic avatar of this renormalisation obstruction.
\end{remark}

\begin{remark}[String field theory interpretation]
\label{rem:sft-bar-identification}
\index{string field theory!bar complex identification}
\index{bar complex!string field theory reading}
Under the bar-BRST identification, the bar differential
$d_{\barB}$ corresponds to the BRST charge~$Q$ of open string
field theory, the bar curvature~$m_0$ to the tadpole, the shadow
amplitude $\mathrm{Sh}_{g,n}(\Theta_\cA)$ to the closed SFT vertex
at genus~$g$ with $n$~punctures, and the MC equation to the
classical closed SFT master equation.
\end{remark}

\section{BV formalism for chiral algebras}

\begin{definition}[BV data for a chiral algebra]
\label{def:bv-data-chiral}
\index{antibracket}
\index{BV algebra!data}
Let $\cA$ be a chiral algebra on a smooth curve~$X$.
The BV formalism associates to~$\cA$:
\begin{enumerate}[label=\textup{(\roman*)}]
\item \emph{Fields}: $\phi \in \mathcal{A}$, sections of the chiral algebra sheaf.
\item \emph{Antifields}: $\phi^+ \in \mathcal{A}^*[1]$, the Serre dual shifted by~$1$.
\item \emph{BV bracket}: $\{\cdot, \cdot\}\colon \mathcal{A} \otimes \mathcal{A}^*[1] \to \mathbb{C}$ of degree~$+1$, an odd Poisson structure.
\item \emph{Action}: $S[\phi, \phi^+]$ satisfying the classical master equation $\{S, S\} = 0$.
\end{enumerate}
\end{definition}

\begin{theorem}[BV complex $=$ geometric bar complex; \ClaimStatusProvedElsewhere{} \cite{CG17}]
\label{thm:bv-bar-geometric}
\index{BV algebra!bar complex identification}
Let $\cA$ be a chiral algebra on $X = \mathbb{P}^1$. The BV complex
$(C_{\mathrm{BV}}(\mathcal{A}), Q_{\mathrm{BV}})$ is isomorphic to the
geometric bar complex:
\[
C_{\mathrm{BV}}(\mathcal{A}) \cong \barB^{\mathrm{ch}}(\mathcal{A}).
\]
Under this isomorphism, $Q_{\mathrm{BV}} = \{S, -\}$ corresponds to the bar
differential $d_{\mathrm{bar}}$. The proof is in Costello--Gwilliam
\cite{CG17}; the following geometric construction explains the
identification.
\end{theorem}

\begin{proof}[Geometric construction supporting \textup{Theorem~\ref{thm:bv-bar-geometric}}]
\emph{Step~1: Bar generators.}

The degree-$n$ bar complex is:
\[\bar{B}^n(\mathcal{A}) = \Omega^*(\overline{C}_{n+1}(X), \mathcal{A}^{\boxtimes (n+1)})\]

The generators are operator insertions $\phi_i \in \mathcal{A}$ at marked points, paired with logarithmic $1$-forms $\eta_{ij} = d\log(z_i - z_j)$ dual to the collision divisors.

\emph{Step~2: BV Bracket.}

The BV bracket is a residue operation on the log-de~Rham complex
of~$\overline{C}_{n+1}(X)$. In local coordinates near the divisor
$D_{ij} = \{z_i = z_j\}$, with $\epsilon = z_i - z_j$:
\[\{\phi(z_i), \eta_{jk}\}
= \delta_{ij}\,\Res_{\epsilon = 0}
\Bigl[\frac{\partial\phi}{\partial z_i}\,
\frac{d\epsilon}{\epsilon(z_i - z_k)}\Bigr]
+ \delta_{ik}\,\Res_{\epsilon = 0}
\Bigl[\frac{\partial\phi}{\partial z_i}\,
\frac{d\epsilon}{\epsilon(z_i - z_j)}\Bigr].\]
Both terms are Poincar\'e residues of logarithmic forms along
normal-crossing divisors of the smooth
variety~$\overline{C}_{n+1}(X)$: the residue map
$\Res_{D_{ij}}\colon \Omega^k_{\log}(\overline{C}_{n+1})
\to \Omega^{k-1}_{D_{ij}}$ is an algebraic morphism of sheaves
on a smooth variety. The factor $1/(z_i - z_k)$ is a pole along
the transverse divisor $D_{ik}$ (Theorem~\ref{thm:FM}(3):
$D_{ij} \cap D_{ik}$ has normal crossings), so on~$D_{ij}$
it restricts to a regular meromorphic function with poles only
along~$D_{ij} \cap D_{ik}$, a codimension-$2$ stratum.
The FM compactification resolves all coincidence singularities
into normal-crossing divisors, and the log-de~Rham complex on
this resolution is the correct algebraic framework: it replaces
distributional regularization with algebraic residues.

The logarithmic $1$-forms $\eta_{ij} = d\log(z_i - z_j)$ generate
$H^1(C_n(X); \mathbb{C})$ by the Arnold--Brieskorn theorem;
for $X = \mathbb{C}$, all $\binom{n}{2}$ classes are linearly independent in degree~$1$
(the Arnold relations constrain products in degree~$\geq 2$), giving
$\dim H^1(C_n(\mathbb{C})) = \binom{n}{2}$. (For a curve of genus~$g > 0$, additional generators from $H^1(X)$ contribute $2gn$ classes.)

\emph{Step~3: Master Equation.}

The classical master equation $\{S, S\} = 0$ is equivalent to $\dzero^2 = 0$ for the
genus-$0$ bar differential:
\[d = d_{\text{strat}} + d_{\text{int}} + d_{\text{res}}\]
\end{proof}

\subsection{Quantum master equation}

\begin{definition}[BV Laplacian]
\label{def:bv-laplacian}
\index{BV Laplacian|textbf}
The BV Laplacian $\Delta_{\mathrm{BV}}$ is the second-order differential operator
\[
\Delta_{\mathrm{BV}} = \sum_i \frac{\partial^2}{\partial \phi_i \partial \phi_i^+},
\]
contracting each field-antifield pair $(\phi_i, \phi_i^+)$.
In the geometric realization on $\barB^{\mathrm{ch}}(\cA)$, the heuristic formula
\[
\Delta_{\mathrm{BV}} \approx \sum_{i<j} \int \delta(z_i - z_j)\,
\frac{\partial}{\partial \eta_{ij}}
\]
inserts a diagonal residue, contracting the logarithmic form $\eta_{ij}$ via a
delta-function along $D_{ij}$. This is the loop-insertion operator in the bar
complex, distinct from the cobar functor $\Omega$ (which recovers $\cA$ from
$B(\cA)$ by bar-cobar inversion, Theorem~\ref{thm:bar-cobar-isomorphism-main}).
The precise algebraic version of this operator is conditional
\textup{(}Theorem~\ref{thm:config-space-bv}\textup{)}.
\end{definition}

\begin{remark}[Heuristic QME/MC comparison; \ClaimStatusHeuristic]
\label{rem:qme-bar-cobar}
\index{quantum master equation|textbf}
At genus~$0$, the classical part of the BV master equation matches the
square-zero bar differential after transporting through the external
comparison of Theorem~\ref{thm:bv-bar-geometric}. Extending this to a
full identification of the quantum master equation
\[
\hbar\,\Delta_{\mathrm{BV}}S + \tfrac{1}{2}\{S,S\}=0
\]
with a Maurer--Cartan equation in a BV-enhanced bar/cobar deformation
complex would require the conditional BV package recorded in
Theorems~\ref{thm:config-space-bv} and~\ref{thm:bv-functor}. That
full comparison is not proved here.

The classical genus-$0$ comparison combines the external BV/bar
identification with the square-zero bar differential on the
uncurved locus. The quantum, Laplacian, and functorial parts are
isolated as conditional content later in this chapter.
\end{remark}

\label{subsec:bv-bar-identification}
\index{BV algebra!bar complex identification}

\begin{remark}[Heuristic genus-\texorpdfstring{$0$}{0} amplitude/bar comparison;
\ClaimStatusHeuristic]\label{rem:genus0-amplitude-bar}
\index{string amplitude!bar complex}
Let $\cA$ be a chiral algebra on $\mathbb{P}^1$, and let
$\omega_1, \ldots, \omega_n$ be cocycles in
$\bar{B}^1(\cA)$. The heuristic genus-$0$ string-amplitude pairing
\[
\mathcal{A}_0(\omega_1, \ldots, \omega_n)
\;=\;
\int_{\overline{\mathcal{M}}_{0,n}}
 \omega_1 \wedge \cdots \wedge \omega_n
\]
parallels the bar-complex Euler characteristic
$\chi(\bar{B}^n_0(\cA))$, where $\bar{B}^n_0$ denotes the
genus-$0$ component of the degree-$n$ bar complex. The exact
identification between these two quantities is not proved in the
manuscript.
The bar complex is realized on
$\overline{C}_n(\mathbb{P}^1) = \overline{\mathcal{M}}_{0,n+1}$
(Theorem~\ref{thm:bar-cobar-isomorphism-main}),
with the bar differential given by residues of the FM propagator
$\eta_{ij} = d\log(z_i - z_j)$ along collision divisors.
The displayed comparison is the expected physics reading of these
configuration-space forms, but the chapter does not prove the required
identification between the physical integration pairing and the Euler
characteristic of the genus-$0$ bar complex.
\end{remark}

\begin{remark}[Genus-$1$ partition function]\label{rem:genus1-bv}
At genus~$1$, the $\hat{A}$-genus gives $F_1 = \kappa(\cA)/24$
\textup{(}Theorem~D\textup{)}. For $\cA = \mathrm{Vir}_c$ with
$\kappa(\mathrm{Vir}_c) = c/2$: $F_1 = c/48$.
\end{remark}

\begin{remark}[Anomaly = curvature]\label{rem:anomaly-curvature-bv}
\index{anomaly!curvature identification}
By Theorem~\ref{thm:anomaly-koszul}, $d_{\mathrm{bar}}^2 = 0$ for $\cA_{\mathrm{tot}} = \cA_{\mathrm{matter}} \otimes \cA_{\mathrm{ghost}}$ if and only if $\kappa_{\mathrm{tot}} = 0$ (equivalently, $c = 26$ for the bosonic string). When $\kappa_{\mathrm{tot}} \neq 0$, the scalar projection of the universal MC class is $\Theta_\cA^{\min} = \kappa(\cA) \cdot \eta \otimes \Lambda$ (Theorem~\ref{thm:explicit-theta}), curving the bar complex at every genus.
\end{remark}

\section{Gauge fixing and BRST}

\subsection{BRST from BV}

\begin{definition}[BRST operator]
\label{def:brst-operator}
\index{BRST cohomology|textbf}
\index{BRST operator|textbf}
The BRST operator $Q_{\mathrm{BRST}}$ arises from gauge fixing the BV action.
Choose a Lagrangian submanifold $\mathcal{L}$ of the total space of fields and
antifields:
\[
Q_{\mathrm{BRST}} = Q_{\mathrm{BV}}\big|_{\mathcal{L}}.
\]
In the chiral algebra context, the BRST charge decomposes by antighost number
(not ghost number):
\[
Q_{\mathrm{BRST}} = Q_0 + Q_1 + Q_2 + \cdots,
\]
where $Q_k$ has antighost number~$k$ and operator dimension~$k-1$.
The total BRST charge has definite ghost number~$+1$.
\end{definition}

\begin{theorem}[BRST cohomology = physical states; \ClaimStatusProvedElsewhere{} \cite{CG17}]
\label{thm:brst-physical-states}
The cohomology of $Q_{\text{BRST}}$ computes physical on-shell states:
\[H^*(Q_{\text{BRST}}) \cong \mathcal{A}_{\text{phys}}\]
\end{theorem}

\begin{example}[Free \texorpdfstring{$bc$}{bc} ghost system]
The $bc$ system has fields $b(z)$ of weight $\lambda$ and $c(z)$ of weight $1-\lambda$, OPE $b(z)c(w) \sim (z-w)^{-1}$, central charge $c_{bc} = 1 - 3(2\lambda - 1)^2$, and BRST operator $Q = \oint c(z) \bigl(T_{\text{matter}}(z) + \tfrac{1}{2}T_{\text{ghost}}(z)\bigr) dz$.

For the bosonic string, $\lambda = 2$ (so $b$ has weight~$2$ and $c$ has weight~$-1$), giving $c_{\text{ghost}} = 1 - 3(3)^2 = -26$ and ghost stress tensor $T_{\text{ghost}} = -2b\,\partial c - (\partial b)\,c$.

On the relative complex $\ker(b_0) \cap \ker(L_0^{\mathrm{tot}})$,
the BRST nilpotence condition is:
\[
Q_{\mathrm{BRST}}^2 = 0
\iff c_{\mathrm{matter}} + c_{\mathrm{ghost}} = 0,
\quad\text{i.e., } c_{\mathrm{matter}} = 26
\]
(on the full complex, $Q_{\mathrm{BRST}}^2 = \tfrac{c-26}{12}\,c_0$;
see Lemma~\ref{lem:brst-nilpotence}).

In the bar complex, the residue component $d_{\mathrm{res}}$
of the bar differential heuristically corresponds to $Q_{\mathrm{BRST}}$
under the identification of Theorem~\ref{thm:bv-bar-geometric}.
\end{example}

\subsection{Coupling to topological gravity}

\begin{theorem}[Transformation law for configuration log forms on
an affine chart; \ClaimStatusProvedHere]
\label{thm:log-form-ghost-law}
\index{ghost!transformation law}
\index{logarithmic form!ghost transformation}
\index{Arnold relations!configuration space}
Let $U \subset X$ be an affine chart with coordinate~$z$ and let
$w : U \to \C$ be a holomorphic change of coordinate with
$w'(z) \neq 0$ on~$U$. On the configuration space
$C_n(U) = \{(z_1,\dots,z_n) \in U^n : z_i \neq z_j\}$ the logarithmic
forms $\eta_{ij} = d\log(z_i - z_j) \in \Omega^{1}_{\log}(C_n(U))$
transform as
\begin{equation}\label{eq:eta-transformation}
\eta_{ij} \;\mapsto\; \eta_{ij}
+ d\log\!\left(\frac{w(z_i) - w(z_j)}{z_i - z_j}\right),
\end{equation}
where the equality is as logarithmic $1$-forms on $C_n(U)$
under the pullback along the $n$-fold self-map
$w^{(n)} : C_n(U) \to C_n(\C)$,
$(z_1,\dots,z_n) \mapsto (w(z_1),\dots,w(z_n))$,
which is well-defined because $w'$ does not vanish on~$U$, so
$w^{(n)}$ preserves the no-collision locus.
\begin{enumerate}[label=\textup{(\roman*)}]
\item \emph{Diagonal limit (fixing the typo).} The ratio
$(w(z_i) - w(z_j))/(z_i - z_j)$ admits the analytic continuation
$\int_0^1 w'(z_j + t(z_i - z_j))\,dt$ to the diagonal; along the
thick diagonal $z_i \to z_j$ this limit equals $w'(z_j)$
(equivalently $w'(z_i)$ when evaluated on the diagonal, since
$z_i = z_j$ there).
\item \emph{Infinitesimal formula.} For an infinitesimal
diffeomorphism $w(z) = z + \epsilon v(z) + O(\epsilon^2)$ generated
by $v \in \mathrm{Vect}(U)$,
\begin{equation}\label{eq:delta-v-eta}
\delta_v \eta_{ij}
= d\!\left(\frac{v(z_i) - v(z_j)}{z_i - z_j}\right).
\end{equation}
This is a Vect$(U)$ $1$-cocycle valued in $\Omega^{1}(C_n(U))$,
with support on the coincidence divisor~$D_{ij}$.
\item \emph{Arnold--Orlik--Solomon quadratic relations.} The forms
$\eta_{ij}$ satisfy
\begin{equation}\label{eq:arnold-3term}
\eta_{ij} \wedge \eta_{jk} + \eta_{jk} \wedge \eta_{ki}
+ \eta_{ki} \wedge \eta_{ij} = 0
\qquad (i,j,k \text{ distinct})
\end{equation}
by the Arnold identity~\cite{Arnold69} on
$C_n(\C)$, equivalently by the Orlik--Solomon
relation~\cite{OrlSol80} of the braid
arrangement~$\cA_n$.
\end{enumerate}
\end{theorem}

\begin{proof}
\emph{Equation \eqref{eq:eta-transformation}.} Factor
$w(z_i) - w(z_j) = (z_i - z_j) \cdot q(z_i, z_j)$ where
$q(z_i, z_j) = \int_0^1 w'(z_j + t(z_i - z_j))\,dt$ is holomorphic on
$U \times U$ with $q|_{\Delta} = w'$; then
$d\log(w(z_i) - w(z_j)) = d\log(z_i - z_j) + d\log q(z_i, z_j)$,
giving \eqref{eq:eta-transformation}.

\emph{Diagonal limit (i).} By the mean-value formula above,
$\lim_{z_i \to z_j} q(z_i,z_j) = w'(z_j)$. The symmetric phrasing
$w'(z_i) = w'(z_j)$ on the diagonal $z_i = z_j$ eliminates the
asymmetry ambiguity.

\emph{Infinitesimal formula (ii).} Set $w(z) = z + \epsilon v(z)$;
then $w(z_i) - w(z_j) = (z_i - z_j) + \epsilon(v(z_i) - v(z_j))$, so
$(w(z_i) - w(z_j))/(z_i - z_j) = 1 + \epsilon\,(v(z_i) - v(z_j))/(z_i - z_j)$.
Taking $d\log$ and linearising in~$\epsilon$ yields
\eqref{eq:delta-v-eta}. This corrects the expression
$d(v'(z_i)/(z_i - z_j))$ that appeared in earlier drafts: the
infinitesimal variation is a \emph{two-point} divided difference,
not a single-point derivative.

\emph{Arnold relation (iii).} The identity~\eqref{eq:arnold-3term}
is Arnold's original cohomology identity~\cite{Arnold69}
for the braid arrangement; it is equivalent to the Orlik--Solomon
presentation of~$H^{*}(C_n(\C),\C)$~\cite{OrlSol80}.
On a general affine chart $U$ the identity holds because it is
local at each triple $(z_i, z_j, z_k)$ and each $\eta_{ij}$ is a
pullback of the analogous form on~$\C^2$.
\end{proof}

\begin{remark}[Scope: which \textup{BRST ghost} this
transformation law is, and which it is \emph{not}]
\label{rem:log-form-ghost-scope}
\index{BRST ghost!scope}
Equation~\eqref{eq:delta-v-eta} is the two-point transgression of
a Vect$(U)$ $1$-cocycle; it is \emph{not} the Lie-derivative
transformation of a weight-$(-1)$ BRST ghost field $c(z)$. The
two are related but distinct:
\begin{enumerate}[label=\textup{(\alph*)}]
\item \emph{Lie-derivative of ghost.} Under $v \in \mathrm{Vect}(U)$
a weight-$(-1)$ ghost field $c$ transforms as
$\delta_v c(z) = (\mathcal{L}_v c)(z) = v(z)\,\partial c(z) -
(\partial v)(z)\,c(z)$ (single point). The Chevalley differential
on $C^{*}(\mathrm{Vect}(U))$ is
$d_{\mathrm{Ch}} c = \tfrac{1}{2}[c, c] = c\,\partial c$, expressing
the Jacobi identity of Vect$(U)$ (nilpotence of Chevalley $d^{2} = 0$).
\item \emph{Two-point transgression.} The formula
\eqref{eq:delta-v-eta} lives on $C_n(U) \supset \Delta$ and is
supported on the coincidence divisor; it is the transgression of
Arnold's Vect-cocycle along the pair of points $(z_i, z_j)$, not
the variation of a ghost at a single point.
\end{enumerate}
These are different constructions and should not be conflated.
The equality "$\delta c = c\,\partial c$" is the Chevalley
$2$-cocycle condition on $\mathrm{Vect}(U)$, not a transformation
law for $c$ under a Vect$(U)$-action.
\end{remark}

\begin{lemma}[Kohno--Drinfeld identification;
\ClaimStatusProvedElsewhere{} \cite{Kohno87,Drinfeld90}]
\label{lem:kohno-drinfeld-log}
\index{Kohno--Drinfeld}
\index{pure braid!Malcev algebra}
The algebra generated over $\Q$ by symbols $t_{ij}$ ($1 \leq i < j \leq n$)
modulo the infinitesimal pure-braid (Arnold) relations
\[
[t_{ij}, t_{kl}] = 0 \quad (\{i,j\} \cap \{k,l\} = \emptyset),
\qquad
[t_{ij}, t_{ik} + t_{jk}] = 0 \quad (i,j,k \text{ distinct})
\]
is isomorphic to the Malcev Lie algebra $\mathfrak{t}_n$ of the
pure braid group $P_n$. Under the Kohno--Drinfeld
isomorphism~\cite{Kohno87,Drinfeld90}, the
cohomology class $[\eta_{ij}] \in H^{1}(C_n(\C),\C)$ is dual to
$t_{ij}$. The resulting Drinfeld--Kohno connection
$\nabla_{\KZ} = d - \sum_{i<j} t_{ij}\,\eta_{ij}$ is flat iff the
Arnold relations \eqref{eq:arnold-3term} hold. This is the
algebraic object to which the forms $\eta_{ij}$ belong; it is
\emph{not} the Chevalley complex of Vect$(U)$.
\end{lemma}

\begin{remark}[Earlier draft errors corrected]
\label{rem:log-form-ghost-corrections}
The earlier draft of
Theorem~\ref{thm:log-form-ghost-law} contained four specific errors,
each listed here for the record and resolved in the statement above:
\begin{enumerate}[label=(E\arabic*)]
\item The diagonal limit was asserted to be $w'(z_i)$, which is
asymmetric in $(i,j)$. The symmetric statement is $w'(z_j)$ (or
equivalently $w'(z_i)$ on the diagonal $z_i = z_j$); see
\eqref{eq:eta-transformation}(i).
\item The infinitesimal formula was written
$\delta_v \eta_{ij} = d(v'(z_i)/(z_i - z_j))$, conflating the
derivative $v'$ with the two-point divided difference
$(v(z_i) - v(z_j))/(z_i - z_j)$; the correct formula is
\eqref{eq:delta-v-eta}.
\item The analogy "exactly as BRST ghosts for diffeomorphisms"
conflated the two-point transgression of a Vect$(U)$-cocycle (which
\eqref{eq:delta-v-eta} is) with the single-point Lie-derivative
transformation of a weight-$(-1)$ ghost field $c(z)$ (which it is
not); see Remark~\ref{rem:log-form-ghost-scope}.
\item The quadratic relation was named "the BRST ghost algebra
($Q_{\mathrm{BRST}}(c) = c\,\partial c$)", confusing Chevalley
nilpotence of a scalar ghost with the Arnold
$3$-term cyclic relation on $C_n(\C)$; these are different
quadratic relations on different complexes. The Arnold relation
\eqref{eq:arnold-3term} is what governs $\eta_{ij}$; the correct
algebraic home is the Malcev algebra $\mathfrak{t}_n$ of
Lemma~\ref{lem:kohno-drinfeld-log}.
\end{enumerate}
\end{remark}

\begin{conjecture}[Bar complex \texorpdfstring{$=$}{=} topological BRST;
\ClaimStatusConjectured]
\label{conj:bar-topological-brst}
The transformation law
(Theorem~\ref{thm:log-form-ghost-law}) and the Kohno--Drinfeld
identification (Lemma~\ref{lem:kohno-drinfeld-log}) extend to a
dg algebra quasi-isomorphism
\[
\bar{B}^{\mathrm{ch}}(\mathcal{A}) \;\simeq\;
C^*_{\mathrm{top}\text{-}\mathrm{BRST}}(\mathcal{A}
\otimes U(\mathfrak{t}_\infty))
\]
on an affine chart~$U$, where $\mathfrak{t}_\infty = \varinjlim \mathfrak{t}_n$
is the colimit Malcev algebra of the tower of pure-braid groups and
the right-hand side is the reduced Chevalley complex of
$\mathcal{A} \otimes U(\mathfrak{t}_\infty)$, not of
$\mathrm{Diff}(U)$.
\end{conjecture>

\section{The genus-\texorpdfstring{$0$}{0} BRST-bar chain map}
\label{sec:brst-bar-chain-map}
\index{BRST!bar complex chain map|textbf}

\subsection{Setup: the BRST complex}

Let $\cA$ be a conformal vertex algebra of central charge~$c$ on
$X = \mathbb{P}^1$, and let $\cA_{\mathrm{gh}} = bc$ denote the
$bc$-ghost system of weights $(2, -1)$ with $c_{\mathrm{gh}} = -26$.
The \emph{BRST complex} of~$\cA$ is the vertex algebra
$\cA_{\mathrm{tot}} = \cA \otimes \cA_{\mathrm{gh}}$ equipped with
the BRST differential:
\begin{equation}\label{eq:brst-differential}
Q_{\mathrm{BRST}}
= \oint \! c(z)\Bigl(T_{\cA}(z) + \tfrac{1}{2}\,T_{\mathrm{gh}}(z)\Bigr)\,dz
= \sum_{n \in \mathbb{Z}} \Bigl(
 \sum_m L^{\cA}_m\, c_{n-m}
 + \tfrac{1}{2}\sum_{p+q=n}(p-q)\, {:}b_p\, c_q\, c_{n-p-q}{:}
\Bigr)
\end{equation}
where $L^{\cA}_n$ are the Virasoro modes of~$\cA$ and $b_n, c_n$ are
the ghost modes with $\{b_m, c_n\} = \delta_{m+n,0}$.

\begin{lemma}[BRST nilpotence; \ClaimStatusProvedElsewhere{} \cite{FGZ86}]
\label{lem:brst-nilpotence}
As operators on $\cA_{\mathrm{tot}}$,
\[
Q_{\mathrm{BRST}}^2 = \tfrac{c - 26}{12}\,c_0.
\]
In particular, $Q_{\mathrm{BRST}}^2 = 0$ on the relative complex
$\cA_{\mathrm{tot}}^{\mathrm{rel}} = \ker(b_0) \cap \ker(L_0^{\mathrm{tot}})$
if and only if $c = 26$.
\end{lemma}

\begin{remark}[BRST nilpotence and the bar construction]
\label{rem:brst-nilpotence-periodicity}
\index{nilpotence-periodicity correspondence!BRST instance}
The condition $Q_{\mathrm{BRST}}^2 = 0$ is the BRST instance of the nilpotence-periodicity correspondence (Remark~\ref{rem:nilpotence-periodicity}): $\kappa_{\mathrm{tot}} = 0$ ensures $d_{\mathrm{bar}}^2 = 0$ on the genus-$0$ bar complex, so the bar differential is a genuine differential. When $\kappa_{\mathrm{tot}} \neq 0$, the curvature term $Q_{\mathrm{BRST}}^2 = \frac{c-26}{12}\,c_0 \neq 0$ obstructs nilpotence.
\end{remark}

The \emph{semi-infinite cohomology} $H^{\infty/2+\bullet}(\cA)$
is the cohomology of $(\cA_{\mathrm{tot}}, Q_{\mathrm{BRST}})$, graded
by ghost number. For the bosonic string ($\cA = \mathcal{H}^{26}$,
$c = 26$), this computes the physical state space.

\subsection{The chain map}

\begin{theorem}[Genus-\texorpdfstring{$0$}{0} BRST-bar quasi-isomorphism;
\ClaimStatusProvedHere]
\label{thm:brst-bar-genus0}
\index{BRST!bar complex quasi-isomorphism}
Let $\cA$ be a conformal vertex algebra on $\mathbb{P}^1$ with
$c(\cA) = 26$ \textup{(}so that
$\cA_{\mathrm{tot}} = \cA \otimes bc$ satisfies
$\kappa_{\mathrm{tot}} = 0$\textup{)}. There exists a chain map
\[
\Phi \colon
 \bigl(\cA_{\mathrm{tot}}^{\mathrm{rel}},\, Q_{\mathrm{BRST}}\bigr)
 \longrightarrow
 \bigl(\barB^{\mathrm{ch}}_0(\cA_{\mathrm{tot}}),\, d_{\mathrm{bar}}\bigr)
\]
from the relative BRST complex to the genus-$0$ bar complex, which is a
quasi-isomorphism of cochain complexes.
\end{theorem}

\begin{proof}
The proof uses a PBW filtration on both complexes, reducing the
comparison to the classical (finite-dimensional) setting where the
result is known.

\emph{Step~1: Conformal weight filtration.}
Both complexes carry a $\mathbb{Z}_{\geq 0}$-filtration by
\emph{total conformal weight}:
\begin{align}
F^{\leq N}\cA_{\mathrm{tot}}^{\mathrm{rel}}
 &= \bigoplus_{\substack{\text{states of total}\\\text{weight } \leq N}}
 \cA_{\mathrm{tot}}^{\mathrm{rel}}, \\
F^{\leq N}\barB^{\mathrm{ch}}_0(\cA_{\mathrm{tot}})
 &= \bigoplus_{\substack{n \geq 0}}
 \Omega^{\leq N}\!\bigl(
 \overline{C}_{n+1}(\mathbb{P}^1),\,
 \cA_{\mathrm{tot}}^{\boxtimes(n+1)}
 \bigr),
\end{align}
where the filtration on bar generators counts the sum of conformal
weights of inserted operators plus the form degree on the
configuration space. Both differentials respect this filtration:
$Q_{\mathrm{BRST}}$ preserves conformal weight (the BRST current
has weight~$1$ and the contour integral lowers by~$1$), and the bar
differential $d_{\mathrm{bar}}$ is computed by OPE residues at
collision divisors, which preserve the total weight.

\emph{Step~2: Associated graded.}
On the associated graded $\mathrm{gr}\,F^{\bullet}$, both complexes
degenerate to classical constructions. For
$\mathrm{gr}\,\cA_{\mathrm{tot}}^{\mathrm{rel}}$: only the
weight-$0$ part of $Q_{\mathrm{BRST}}$ survives, which is the
Chevalley--Eilenberg differential of the finite-dimensional Lie algebra
$\mathfrak{g}_0 = \cA_{(1)}$ (the weight-$1$ subspace, viewed as a Lie
algebra via the zero-mode bracket $[a, b] = a_{(0)} b$).
For $\mathrm{gr}\,\barB^{\mathrm{ch}}_0(\cA_{\mathrm{tot}})$: only
the leading OPE pole contributes to the bar differential, giving
the classical bar complex $\barB(\mathfrak{g}_0)$ of the
finite-dimensional Lie algebra~$\mathfrak{g}_0$.

In both cases, the associated graded complex is the
Chevalley--Eilenberg / bar complex of a finite-dimensional Lie algebra.
By Arkhipov~\cite{Arkhipov97}, there is a natural quasi-isomorphism
\begin{equation}\label{eq:arkhipov-classical}
\Phi_0 \colon C^{\bullet}_{\mathrm{CE}}(\mathfrak{g}_0)
 \xrightarrow{\;\sim\;}
 \barB(\mathfrak{g}_0)
\end{equation}
given by the bar-CE comparison map: on generators,
$\Phi_0(x_1 \wedge \cdots \wedge x_n)
= [x_1 | \cdots | x_n]_{\mathrm{bar}}$
(the corresponding bar element), extended by the shuffle product
to all degrees. This is a classical result in homological algebra
(e.g., Loday--Vallette~\cite{LV12}, \S2.2): the bar complex and
the Chevalley--Eilenberg complex of a Lie algebra are
quasi-isomorphic via the antisymmetrization map.

\emph{Step~3: Spectral sequence comparison.}
The conformal weight filtration induces convergent spectral sequences:
\begin{align}
{}^{\mathrm{BRST}}E_1^{p,q}
 &= H^{p+q}\bigl(\mathrm{gr}^p\, \cA_{\mathrm{tot}}^{\mathrm{rel}}\bigr)
 \;\Longrightarrow\;
 H^{p+q}(Q_{\mathrm{BRST}}), \\
{}^{\mathrm{bar}}E_1^{p,q}
 &= H^{p+q}\bigl(\mathrm{gr}^p\, \barB^{\mathrm{ch}}_0(\cA_{\mathrm{tot}})\bigr)
 \;\Longrightarrow\;
 H^{p+q}(d_{\mathrm{bar}}).
\end{align}
Both spectral sequences converge because the filtrations are bounded
below (conformal weight $\geq 0$) and exhaustive (every element has
finite weight), and the complexes are complete with respect to the
filtration (the vertex algebra is $\mathbb{Z}_{\geq 0}$-graded by
conformal weight with finite-dimensional graded pieces).

The map $\Phi_0$ of Step~2 defines an isomorphism
${}^{\mathrm{BRST}}E_1^{p,q} \xrightarrow{\sim}
{}^{\mathrm{bar}}E_1^{p,q}$ at the $E_1$ page: at each conformal
weight~$p$, the associated graded complexes are identified by~$\Phi_0$,
which is a quasi-isomorphism.

\emph{Step~4: Lifting to a chain map.}
The isomorphism at the $E_1$ level lifts to a chain map on the full
complexes by the standard spectral sequence comparison theorem
(Eilenberg--Moore~\cite{EM66}; see also
Weibel~\cite[Theorem~5.2.12]{Weibel94}): given two filtered complexes
with a map of $E_1$ pages that is an isomorphism, there exists a chain
map between the unfiltered complexes inducing this isomorphism, and any
such chain map is a quasi-isomorphism.

Concretely, $\Phi$ is constructed inductively on filtration degree.
At level $N = 0$, $\Phi$ is the classical map~$\Phi_0$. At level~$N$,
given $\Phi|_{F^{\leq N-1}}$, one extends to $F^{\leq N}$ by solving
a lifting problem: the obstruction to extending lies in
$H^{n+1}(\mathrm{gr}^N\,\barB^{\mathrm{ch}}_0)$, which is the
$(n+1)$-st cohomology of the classical bar complex at weight~$N$.
Since the classical bar complex of a Lie algebra is acyclic in
sufficiently high degree (by the Koszul resolution), the obstructions
vanish degree by degree, and the induction proceeds.

The map $\Phi$ so constructed is a filtered chain map inducing an
isomorphism on $E_1$ pages. By the Eilenberg--Moore comparison
theorem, $\Phi$ is a quasi-isomorphism.
\end{proof}

\begin{remark}[BV/BRST reading of the master square]
\label{rem:bv-convergence}
\index{BV complex!convergence with bar complex}
Theorem~\ref{thm:brst-bar-genus0} realizes the parallel-track identifications of Remarks~\ref{rem:bv-parallel-track} and~\ref{rem:anomaly-parallel-track}: in~\eqref{eq:master-square}, $\barB_X(\cA)$ is the BV complex, $\barB_X(\cA^!)$ the ghost complex, $\mathbb{D}_{\mathrm{Ran}}$ the field--antifield exchange, and $\dfib^{\,2} = \kappa \cdot \omega_1$ the conformal anomaly. The square commutes precisely when $\kappa = 0$.
\end{remark}

\begin{remark}[Explicit form of \texorpdfstring{$\Phi$}{math}]\label{rem:brst-bar-explicit}
For the bosonic string $\cA = \mathcal{H}^{26}$, the chain map
$\Phi$ can be described concretely in low degrees:
\begin{itemize}
\item \emph{Degree~$0$}: $\Phi(|0\rangle) = 1 \in \barB^0$.
 A BRST-closed state of ghost number~$0$ maps to the unit in
 the bar complex.
\item \emph{Degree~$1$}: $\Phi(c_1 \alpha_{-1}^i |0\rangle)
 = [\alpha^i] \in \barB^1$, where $\alpha^i$ are the Heisenberg
 generators. The physical vertex operators (momentum eigenstates)
 map to bar generators.
\item \emph{Degree~$2$}: $\Phi(c_0 c_1 |0\rangle) = \eta_{12}$,
 where $\eta_{12} = d\log(z_1 - z_2)$ is the FM propagator.
 The ghost bilinear maps to the logarithmic form on the
 configuration space. This is the ghost--propagator identification.
\end{itemize}
In general, ghost insertions $c_{n_1} \cdots c_{n_k}$ map to
products of logarithmic forms $\eta_{i_1 j_1} \wedge \cdots \wedge
\eta_{i_r j_r}$ on the configuration space, with the precise
correspondence determined by the residue map of the FM propagator.
\end{remark}

\begin{corollary}[Physical anomaly cancellation at genus \texorpdfstring{$0$}{0};
\ClaimStatusProvedHere]
\label{cor:anomaly-physical-genus0}
\index{anomaly cancellation!BRST proof}
The mathematical condition $\kappa_{\mathrm{tot}} = 0$
\textup{(}Theorem~\textup{\ref{thm:anomaly-koszul})} is equivalent,
at genus~$0$, to the BRST anomaly cancellation
$c_{\mathrm{matter}} + c_{\mathrm{ghost}} = 0$. Specifically:
\begin{enumerate}[label=\textup{(\roman*)}]
\item $Q_{\mathrm{BRST}}^2 = 0$ on $\cA_{\mathrm{tot}}^{\mathrm{rel}}$
 if and only if $c(\cA) = 26$ \textup{(}Lemma~\textup{\ref{lem:brst-nilpotence})}.
\item $d_{\mathrm{bar}}^2 = 0$ on $\barB^{\mathrm{ch}}_0(\cA_{\mathrm{tot}})$
 if and only if $\kappa_{\mathrm{tot}} = 0$, i.e., $c(\cA) = 26$
 \textup{(}Theorem~\textup{\ref{thm:anomaly-koszul})}.
\item The chain map $\Phi$ of Theorem~\textup{\ref{thm:brst-bar-genus0}}
 identifies both anomaly cancellation conditions: both sides are well-defined
 cochain complexes precisely when $c = 26$, and $\Phi$ is a quasi-isomorphism
 relating their cohomologies.
\end{enumerate}
\end{corollary}

\begin{proof}
Parts (i) and (ii) are the cited results. Part (iii):
Theorem~\ref{thm:brst-bar-genus0} constructs $\Phi$ under
the hypothesis $c = 26$, giving a quasi-isomorphism that
identifies $H^*(Q_{\mathrm{BRST}}) \cong H^*(d_{\mathrm{bar}})$.
If $c \neq 26$, the BRST side has $Q^2 = \frac{c-26}{12}\,c_0 \neq 0$
and the bar side has $\dfib^{\,2} = \kappa_{\mathrm{tot}} \cdot \omega_1
\cdot \mathrm{id} \neq 0$: neither is a cochain complex,
and both obstructions are controlled by the same numerical invariant
$c - 26$. The dictionary
\[
\kappa_{\mathrm{tot}}
= \kappa(\cA_{\mathrm{matter}}) + \kappa(\cA_{\mathrm{ghost}})
= \tfrac{c}{2} + \tfrac{-26}{2}
= \tfrac{c-26}{2}
\]
(Computation~\ref{comp:virasoro-curvature} for $\kappa(\mathrm{Vir}_c) = c/2$;
$\kappa(\text{$bc$-ghost}, \lambda{=}2) = c_{\mathrm{gh}}/2 = -13$)
converts the BRST obstruction $Q^2 \propto (c-26)$ into
the bar obstruction $\dfib^{\,2} \propto \kappa_{\mathrm{tot}}$.
\end{proof}

\begin{remark}[Attribution: quantum complementarity is a Volume~I theorem]
\label{rem:quantum-complementarity-vol1-attribution}
All invocations of Theorem~\ref{V1-thm:quantum-complementarity-main} in
this chapter refer to the Lagrangian-polarization complementarity
theorem inscribed in Volume~I, chapter on higher-genus
complementarity: for a chiral Koszul pair $(\cA,\cA^!)$ on a smooth
projective curve~$X$, the ambient complex $\mathbf{C}_g(\cA)$ at each
genus~$g$ decomposes as $\mathbf{Q}_g(\cA) \oplus \mathbf{Q}_g(\cA^!)$,
and Verdier duality makes each summand Lagrangian for the induced
pairing of cohomological degree $-(3g{-}3)$; at the cohomological
shadow level this gives
$H^\ast(\overline{\cM}_g, \cZ(\cA)) = Q_g(\cA) \oplus Q_g(\cA^!)$.
The family-dependent scalar complementarity
$\kappa(\cA) + \kappa(\cA^!) \in \{0, 13, \tfrac{250}{3}, \tfrac{98}{3},
\ldots\}$ (free/Kac--Moody, Virasoro, $\cW_3$, Bershadsky--Polyakov)
is the $g=0$ numerical shadow. The present chapter uses Volume~I's
theorem as imported input; no independent proof is attempted here.
\end{remark}

\begin{proposition}[Koszul duality preserves BRST anomaly cancellation; \ClaimStatusProvedHere]
\label{prop:koszul-brst-anomaly-preservation}
\index{Koszul duality!BRST anomaly preservation}
\index{anomaly cancellation!Koszul duality}
Let $\cA_{\mathrm{tot}} = \cA_{\mathrm{matter}} \otimes \cA_{\mathrm{ghost}}$
be a BRST-quantized gauge system with total modular characteristic
$\kappa_{\mathrm{tot}} = \kappa(\cA_{\mathrm{matter}}) +
\kappa(\cA_{\mathrm{ghost}}) = 0$ at the classical level. Assume the matter
sector $\cA_{\mathrm{matter}}$ lies in the free (Heisenberg) or Kac--Moody
family, so that the complementarity sum $\kappa(\cA_{\mathrm{matter}})
+ \kappa(\cA_{\mathrm{matter}}^!) = 0$ holds
\textup{(}Theorem~\ref{V1-thm:quantum-complementarity-main}, free\slash Kac--Moody
branch\textup{)}. Let $\widetilde{\cA}_{\mathrm{ghost}}$ be a ghost sector
obtained by transporting $\cA_{\mathrm{ghost}}$ through a matching
Koszul involution: an involution of the ghost system whose modular
characteristic satisfies
$\kappa(\widetilde{\cA}_{\mathrm{ghost}}) = -\kappa(\cA_{\mathrm{ghost}})$.
Then the dual total system
\[
 \cA_{\mathrm{tot}}^{\vee}
 \;=\; \cA_{\mathrm{matter}}^! \otimes \widetilde{\cA}_{\mathrm{ghost}}
\]
has vanishing total characteristic
\[
 \kappa(\cA_{\mathrm{tot}}^{\vee})
 \;=\; \kappa(\cA_{\mathrm{matter}}^!)
 + \kappa(\widetilde{\cA}_{\mathrm{ghost}})
 \;=\; -\kappa(\cA_{\mathrm{matter}})
 - \kappa(\cA_{\mathrm{ghost}})
 \;=\; -\kappa_{\mathrm{tot}}
 \;=\; 0,
\]
and the BRST anomaly cancellation of
Corollary~\textup{\ref{cor:anomaly-physical-genus0}} holds for
$\cA_{\mathrm{tot}}^{\vee}$ at genus~$0$.
\end{proposition}

\begin{proof}
By Theorem~D \textup{(}$\kappa$-additivity under tensor product of chirally
Koszul algebras, uniform-weight lane; see
Theorem~\ref{thm:explicit-theta}\textup{)},
$\kappa(\cA_{\mathrm{matter}}^! \otimes \widetilde{\cA}_{\mathrm{ghost}})
= \kappa(\cA_{\mathrm{matter}}^!) + \kappa(\widetilde{\cA}_{\mathrm{ghost}})$.
The free\slash Kac--Moody complementarity sum gives
$\kappa(\cA_{\mathrm{matter}}^!) = -\kappa(\cA_{\mathrm{matter}})$,
and the matching involution hypothesis gives
$\kappa(\widetilde{\cA}_{\mathrm{ghost}})
= -\kappa(\cA_{\mathrm{ghost}})$, hence the total is
$-\kappa_{\mathrm{tot}} = 0$. The genus-$0$ BRST cancellation then
follows from Corollary~\ref{cor:anomaly-physical-genus0} applied to
$\cA_{\mathrm{tot}}^{\vee}$.
\end{proof}

\begin{remark}[Scoping: why the ghost involution matters]
\label{rem:koszul-brst-scoping}
The hypothesis on $\widetilde{\cA}_{\mathrm{ghost}}$ is essential and
cannot be dropped. Consider the bosonic string:
$\kappa(\cA_{\mathrm{matter}}) = 13$ ($26$ free scalars: each scalar
contributes $\kappa = 1/2$, so the total is $13$);
$\kappa(\cA_{\mathrm{ghost}}) = -13$ ($bc$ system at $\lambda = 2$,
$c_{bc} = 1 - 3(2{\cdot}2{-}1)^2 = -26$, $\kappa_{bc} = c_{bc}/2 = -13$);
total $\kappa_{\mathrm{tot}} = 13 + (-13) = 0$.
Applying Koszul duality to the matter sector alone sends
$\kappa(\cA_{\mathrm{matter}}) \mapsto -13$
\textup{(}complementarity for the free\slash Heisenberg family,
$\kappa(\mathcal{H}_\kappa) + \kappa(\mathcal{H}_\kappa^!) = 0$;
this is the free branch of
Theorem~\textup{\ref{V1-thm:quantum-complementarity-main}}, \emph{not}
the Virasoro branch where $\kappa + \kappa' = 13$\textup{)}.
If one keeps the \emph{same} ghost sector $\cA_{\mathrm{ghost}}$
unchanged, the new total is $-13 + (-13) = -26 \neq 0$, and BRST
anomaly cancellation fails. The proposition therefore requires that
the ghost sector also be transported by a Koszul-compatible involution
that flips $\kappa_{\mathrm{ghost}}$. For the bosonic string, the
matching involution on the $bc$ system is the one sending
$(\lambda, 1 - \lambda) \mapsto (-1, 2)$, which reverses the sign of
$c_{bc}$ at the normalization ambiguity of the $bc$ system
\textup{(}Remark~\textup{\ref{rem:ghost-superghost-koszul})}.
Without such a transport, matter Koszul duality alone does not
preserve the anomaly.
\end{remark}

\begin{computation}[Superstring ghost/superghost Koszul dual;
\ClaimStatusProvedHere]
\label{comp:superstring-ghost-koszul}
\index{superstring!ghost Koszul dual}
\index{$\beta\gamma$ system!superghost Koszul dual}
The bosonic string ghost system is the $bc$ pair at $\lambda = 2$
with $c_{bc} = 1 - 3(2{\cdot}2 - 1)^2 = -26$ and
$\kappa_{bc} = c_{bc}/2 = -13$. The complementarity relation
$\kappa + \kappa' = -26 \neq 13$ shows that the $bc$ system is
\emph{not} self-dual under the Virasoro self-duality locus
$c = 13$ \textup{(}Virasoro self-dual at
$\kappa + \kappa' = 13$\textup{)}: the $bc$ ghost is self-dual
under a different involution \textup{(}the one described in
Remark~\textup{\ref{rem:koszul-brst-scoping})} that fixes
the bilinear pairing but reverses the total $\kappa$.

For the superstring, the ghost sector enlarges to
$(bc) \otimes (\beta\gamma)$, where the superghost
$\beta\gamma$ system sits at $\lambda = 3/2$. By the
formula $c_{\beta\gamma}(\lambda) = 2(6\lambda^2 - 6\lambda + 1)$ \[
 c_{\beta\gamma}(\tfrac{3}{2})
 \;=\; 2\bigl(6 \cdot \tfrac{9}{4}
 - 6 \cdot \tfrac{3}{2} + 1\bigr)
 \;=\; 2\bigl(\tfrac{27}{2} - 9 + 1\bigr)
 \;=\; 2 \cdot \tfrac{11}{2}
 \;=\; 11,
\]
hence $\kappa_{\beta\gamma} = c_{\beta\gamma}/2 = 11/2$, and the
Koszul dual satisfies $\kappa'_{\beta\gamma} = -11/2$.

The full superghost characteristic is
\[
 \kappa_{\mathrm{ghost}}^{\mathrm{super}}
 \;=\; \kappa_{bc} + \kappa_{\beta\gamma}
 \;=\; -13 + \tfrac{11}{2}
 \;=\; -\tfrac{15}{2}.
\]
Superstring BRST anomaly cancellation requires
$\kappa_{\mathrm{tot}} = \kappa_{\mathrm{matter}}
+ \kappa_{\mathrm{ghost}}^{\mathrm{super}} = 0$,
hence $\kappa_{\mathrm{matter}} = 15/2$, equivalently
$c_{\mathrm{matter}} = 15$. For the $N = 1$ superconformal
matter of $d$ flat spacetime dimensions
($d$ free scalars contributing $d$ to $c$, plus $d$ free
Majorana fermions contributing $d/2$), this gives
$c_{\mathrm{matter}} = d + d/2 = 3d/2 = 15$, hence
$d = 10$: the critical dimension of the superstring. The
complementarity sum for the superghost is
\[
 \kappa_{\mathrm{ghost}}^{\mathrm{super}}
 + (\kappa_{\mathrm{ghost}}^{\mathrm{super}})'
 \;=\; -\tfrac{15}{2} + \tfrac{15}{2}
 \;=\; 0,
\]
confirming that the superstring BRST anomaly cancellation
\textup{(}at $d = 10$\textup{)} is preserved by the matching
Koszul involution of Proposition~\ref{prop:koszul-brst-anomaly-preservation}.
\end{computation}

\begin{remark}[The ghost--superghost Koszul involution;
\ClaimStatusProvedElsewhere]
\label{rem:ghost-superghost-koszul}
\index{ghost!Koszul involution}
\index{superghost!Koszul involution}
The preceding computation identifies the ghost characteristic in
each of the two string backgrounds:
\begin{itemize}[nosep]
\item bosonic string: ghost $= bc$ at $\lambda = 2$, with
 $\kappa = -13$;
\item superstring: ghost $= bc$ at $\lambda = 2$ plus $\beta\gamma$
 at $\lambda = 3/2$, with total $\kappa = -15/2$.
\end{itemize}
In each case, the Koszul dual is obtained by a single involution
that reverses the sign of the ghost characteristic while preserving
the bilinear pairing. For the $bc$ system, the involution is
the one identified by Friedan--Martinec--Shenker~\cite{FMS86}
relating the $bc$ and $\beta\gamma$ systems at bosonised level;
for the superstring combined ghost sector, the involution acts
diagonally on both factors. The statement
$\kappa_{\mathrm{ghost}}^{\mathrm{bos}} = -13$ and
$\kappa_{\mathrm{ghost}}^{\mathrm{super}} = -15/2$ are the
numerical invariants of these involutions, matching the standard
critical dimensions $d = 26$ and $d = 10$ of
Polchinski~\cite{Polchinski1998}. The involution is
\emph{not} the Virasoro self-duality at $c = 13$
\textup{(}which would require $\kappa + \kappa' = 13$\textup{)}:
the ghost involution is native to the $bc$\slash$\beta\gamma$
family and operates at a different fixed locus.
\end{remark}

\begin{remark}[Higher genus]
\label{rem:brst-bar-higher-genus}
At genus $g \geq 1$, the BRST complex requires additional data:
the path integral measure on $\overline{\mathcal{M}}_g$ and the
Costello renormalization framework~\cite{costello-renormalization} for handling
ultraviolet divergences. The bar complex side is well-defined at
all genera via configuration spaces on $\Sigma_g$
(Chapter~\ref{chap:higher-genus}). The chain map $\Phi$ should
extend to a map of genus-graded complexes, with the genus-$g$
component involving Costello's counterterm expansion. This
extension is the subject of Programme~VI-a in
\S\ref{sec:modular-koszul-programme}.
\end{remark}

\section{Semi-infinite cohomology and the bar complex}
\label{sec:semi-infinite-bar}
\index{semi-infinite cohomology|textbf}

\subsection{Semi-infinite cohomology}

\begin{definition}[Semi-infinite cohomology {\cite{FGZ86}}]
\label{def:semi-infinite-cohomology}
\index{semi-infinite cohomology!definition}
Let $\widehat{L} = \bigoplus_{n \in \mathbb{Z}} L_n$ be a
$\mathbb{Z}$-graded Lie algebra with $\dim L_n < \infty$, and let
$M$ be a $\widehat{L}$-module. Write
$\widehat{L} = L_- \oplus L_0 \oplus L_+$ for the decomposition
into negative, zero, and positive modes. The \emph{semi-infinite
complex} is:
\[
C^{\infty/2+n}(\widehat{L}, M)
= M \otimes \textstyle\bigwedge\nolimits^{\infty/2+n}\! \widehat{L}^*
\]
where $\bigwedge^{\infty/2+n}$ denotes the semi-infinite wedge
power: the ``Dirac sea'' vacuum is the semi-infinite form
$\xi_0^* \wedge \xi_{-1}^* \wedge \xi_{-2}^* \wedge \cdots$
(wedging all negative-mode duals), and degree $n$ is measured
relative to this vacuum. The semi-infinite differential
$Q_{\mathrm{si}}$ combines the Chevalley--Eilenberg differential
with a ``charge anomaly'' correction term that ensures
$Q_{\mathrm{si}}^2 = 0$ without any restriction on the central charge.
The \emph{semi-infinite cohomology} is
$H^{\infty/2+\bullet}(\widehat{L}, M) =
H^*(C^{\infty/2+\bullet}(\widehat{L}, M), Q_{\mathrm{si}})$.
\end{definition}

\begin{remark}[Semi-infinite vs.\ string BRST]
\label{rem:semi-infinite-vs-brst}
The string BRST complex of \S\ref{sec:brst-bar-chain-map} couples
$\cA$ to $bc$-ghosts and requires $c = 26$ for $Q^2 = 0$
(Lemma~\ref{lem:brst-nilpotence}). Semi-infinite cohomology is
different: it is defined for any $\mathbb{Z}$-graded Lie algebra
$\widehat{L}$ and any module $M$, with $Q_{\mathrm{si}}^2 = 0$
unconditionally. The ghosts are ``built in'' to the semi-infinite
wedge product rather than introduced as an external $bc$-system.
For the string ($\widehat{L} = \mathrm{Vir}$, $M = \cA$),
semi-infinite cohomology of the Virasoro algebra \emph{does}
reduce to string BRST cohomology when $c = 26$~\cite{FGZ86}.
\end{remark}

\subsection{The bar--semi-infinite identification for Kac--Moody algebras}

\begin{theorem}[Bar complex = semi-infinite complex for KM;
\ClaimStatusProvedHere]
\label{thm:bar-semi-infinite-km}
\index{semi-infinite cohomology!bar complex identification}
\index{BRST!semi-infinite identification}
Let $\widehat{\fg}_k$ be the universal affine vertex algebra at level
$k \neq -h^\vee$. There is a quasi-isomorphism of filtered cochain
complexes:
\begin{equation}\label{eq:bar-semi-infinite}
\Psi \colon
 \bigl(\barB^{\mathrm{ch}}(\widehat{\fg}_k),\, d_{\mathrm{bar}}\bigr)
 \xrightarrow{\;\sim\;}
 \bigl(C^{\infty/2+\bullet}(\widehat{\fg}_k, V_k),\, Q_{\mathrm{si}}\bigr)
\end{equation}
identifying the chiral bar complex with the semi-infinite complex
\textup{(}with coefficients in the vacuum module $V_k$\textup{)}.
In particular:
\begin{equation}\label{eq:bar-equals-semi-infinite}
H^*(\barB^{\mathrm{ch}}(\widehat{\fg}_k))
\;\cong\;
H^{\infty/2+*}(\widehat{\fg}_k,\, V_k).
\end{equation}
\end{theorem}

\begin{proof}
The proof parallels that of
Theorem~\ref{thm:brst-bar-genus0}, replacing the string BRST
complex with the semi-infinite complex and removing the
$c = 26$ hypothesis.

\emph{Step~1: Conformal weight filtration.}
Both complexes carry bounded-below, exhaustive
$\mathbb{Z}_{\geq 0}$-filtrations by total conformal weight.
On the bar side, this is the PBW filtration of
Theorem~\ref{thm:brst-bar-genus0}, Step~1. On the semi-infinite
side, the conformal weight filtration arises from the
$L_0$-grading on $V_k$ and the mode grading on
$\bigwedge^{\infty/2+\bullet}\!\widehat{\fg}^*_k$. Both
differentials ($d_{\mathrm{bar}}$ and $Q_{\mathrm{si}}$) respect
these filtrations: $d_{\mathrm{bar}}$ preserves total conformal
weight (OPE residues preserve weight), and $Q_{\mathrm{si}}$
preserves it by construction~\cite{FGZ86}.

\emph{Step~2: Associated graded.}
On the associated graded $\mathrm{gr}\,F^\bullet$,
$\mathrm{gr}\,\barB^{\mathrm{ch}}(\widehat{\fg}_k)$ reduces
to the classical bar complex $\barB(\fg)$ of the
finite-dimensional Lie algebra $\fg$
(only the leading OPE pole survives, giving the Lie bracket),
while $\mathrm{gr}\,C^{\infty/2+\bullet}(\widehat{\fg}_k, V_k)$
reduces to the Chevalley--Eilenberg complex
$C^*_{\mathrm{CE}}(\fg, \mathbb{C})$
(the semi-infinite vacuum collapses to $\bigwedge^*\!\fg^*$
when mode corrections are dropped;
see~\cite{FGZ86}, Proposition~2.1).
Both associated graded complexes compute the Lie algebra cohomology
$H^*(\fg, \mathbb{C})$. The classical bar--CE comparison map
$\Phi_0$ of Arkhipov~\cite{Arkhipov97}
(cf.\ \eqref{eq:arkhipov-classical}) provides a natural
quasi-isomorphism $\barB(\fg) \xrightarrow{\sim}
C^*_{\mathrm{CE}}(\fg, \mathbb{C})$.

\emph{Step~3: Spectral sequence comparison.}
The conformal weight filtrations induce convergent spectral sequences
on both sides (bounded below, exhaustive, finite-dimensional
graded pieces). The map $\Phi_0$ induces an isomorphism at the
$E_1$ page. By the Eilenberg--Moore comparison theorem
(\cite{EM66}; cf.\ \cite[Theorem~5.2.12]{Weibel94}), $\Phi_0$ lifts
to a filtered chain map
$\Psi \colon \barB^{\mathrm{ch}}(\widehat{\fg}_k) \to
C^{\infty/2+\bullet}(\widehat{\fg}_k, V_k)$
that is a quasi-isomorphism.

\emph{Step~4: No anomaly obstruction.}
Unlike Theorem~\ref{thm:brst-bar-genus0}, no central charge
condition is needed. On the bar side, $d_{\mathrm{bar}}^2 = 0$
by the Arnold relations
(Theorem~\ref{thm:arnold-three}); the curvature
$m_0 = \kappa \cdot \mathbf{1}$ contributes to the $A_\infty$
structure but does not obstruct the bar differential. On the
semi-infinite side, $Q_{\mathrm{si}}^2 = 0$ by
construction~\cite{FGZ86}: the semi-infinite wedge product
absorbs the charge anomaly that would otherwise require
$c = 26$. Both differentials square to zero at every level
$k \neq -h^\vee$, and the spectral sequence comparison
applies without restriction.

The condition $k \neq -h^\vee$ ensures that the Sugawara
construction is well-defined, providing the $L_0$-grading
that drives the filtration. At critical level
$k = -h^\vee$, the Sugawara stress tensor is undefined, but the
semi-infinite complex is still defined using the loop rotation
grading; the identification at critical level follows from the
Feigin--Frenkel theorem~\cite{Feigin-Frenkel}
(cf.\ Theorem~\ref{thm:critical-level-structure}).
\end{proof}

\begin{remark}[String BRST as a special case of semi-infinite cohomology]
\label{rem:string-brst-special-case}
Theorem~\ref{thm:bar-semi-infinite-km} specializes to recover
Theorem~\ref{thm:brst-bar-genus0}: take $\widehat{\fg}_k = \mathrm{Vir}$
and $\cA_{\mathrm{tot}} = \mathcal{H}^{26} \otimes bc$ with
$c_{\mathrm{tot}} = 0$. By~\cite{FGZ86}, the string BRST complex is the
semi-infinite complex of the Virasoro algebra acting on~$\mathcal{H}^{26}$;
the chain map~$\Psi$ of Theorem~\ref{thm:bar-semi-infinite-km} then recovers
the map~$\Phi$ of Theorem~\ref{thm:brst-bar-genus0}. At $c = 26$, both
$Q_{\mathrm{BRST}}^2 = \frac{c-26}{12}\,c_0$ and $Q_{\mathrm{si}}^2$
vanish, confirming the coincidence; away from $c = 26$, the
$\frac{c-26}{12}\,c_0$ term measures the mismatch.
\end{remark}

\subsection{Anomaly duality from complementarity}

\begin{corollary}[Anomaly duality for Kac--Moody pairs;
\ClaimStatusProvedHere]
\label{cor:anomaly-duality-km}
\index{anomaly cancellation!Kac--Moody duality}
\index{Feigin--Frenkel duality!anomaly interpretation}
Let $\widehat{\fg}_k$ be an affine Kac--Moody algebra at generic
level $k$, with Feigin--Frenkel dual
$\widehat{\fg}_{k'} = \widehat{\fg}_{-k-2h^\vee}$. Then:
\begin{enumerate}[label=\textup{(\alph*)}]
\item \emph{Bar cohomology = semi-infinite cohomology.}
The bar cohomology computations of
Part~\ref{part:characteristic-datum} \textup{(}Riordan numbers for
$\widehat{\mathfrak{sl}}_2$, the Hilbert series of
Theorem~\textup{\ref{thm:bar-cohomology-level-independence})}
are semi-infinite cohomology computations:
\[
\dim H^n(\barB^{\mathrm{ch}}(\widehat{\fg}_k))
= \dim H^{\infty/2+n}(\widehat{\fg}_k, V_k).
\]

\item \emph{Curvature = conformal anomaly.}
The genus-$1$ curvature coefficient
$\kappa(\widehat{\fg}_k) = (k+h^\vee)\dim\fg/(2h^\vee)$
\textup{(}Theorem~\textup{\ref{thm:genus-universality}(i))} equals
the semi-infinite anomaly coefficient. For affine Kac--Moody pairs,
the Feigin--Frenkel anti-symmetry
$\kappa(\widehat{\fg}_k) + \kappa(\widehat{\fg}_{-k-2h^\vee}) = 0$
of Theorem~\textup{\ref{thm:genus-universality}(ii)} is the
anomaly cancellation for the Feigin--Frenkel pair: the
conformal anomalies of $\widehat{\fg}_k$ and
$\widehat{\fg}_{-k-2h^\vee}$ cancel.
\textup{(}This anti-symmetry holds for KM/free-field families where the
duality is level-negating; it fails for Virasoro, where
$\kappa(\mathrm{Vir}_c) + \kappa(\mathrm{Vir}_{26-c}) = 13 \neq 0$.\textup{)}

\item \emph{Complementarity = anomaly duality.}
The quantum complementarity theorem
\textup{(}Theorem~\textup{\ref{V1-thm:quantum-complementarity-main})}
applied to the Koszul pair
$(\widehat{\fg}_k, \widehat{\fg}_{-k-2h^\vee})$ gives:
\[
Q_g(\widehat{\fg}_k) \oplus Q_g(\widehat{\fg}_{-k-2h^\vee})
= H^*(\overline{\mathcal{M}}_g,\, Z(\widehat{\fg}_k)).
\]
Via the bar--semi-infinite identification, the left side is a
direct sum of semi-infinite cohomology spaces on
$\overline{\mathcal{M}}_g$: the genus-$g$ quantum corrections
of the WZW model at level~$k$ and its Feigin--Frenkel dual sum to
the full moduli space cohomology. This is the mathematical
content of anomaly cancellation for Kac--Moody pairs at all genera.

\item \emph{Universal MC duality.}
The universal Maurer--Cartan classes satisfy
$\Theta_{\widehat{\fg}_k} + \Theta_{\widehat{\fg}_{-k-2h^\vee}}
= 0$ (Theorem~\ref{thm:explicit-theta}), since
$\kappa(\widehat{\fg}_k) + \kappa(\widehat{\fg}_{-k-2h^\vee}) = 0$
by Theorem~\ref{thm:genus-universality}(ii).
The chain-level mechanism underlying~(c) is:
the Verdier involution exchanges
$\Theta_{\widehat{\fg}_k}$ and $-\Theta_{\widehat{\fg}_{k'}}$,
producing the Lagrangian splitting of
Theorem~\ref{V1-thm:quantum-complementarity-main}.
\end{enumerate}
\end{corollary}

\begin{proof}
Part~(a) is the cohomology-level consequence of
Theorem~\ref{thm:bar-semi-infinite-km}. Part~(b): the
curvature $m_0 = \kappa \cdot \mathbf{1}$ in the bar complex
corresponds, under the quasi-isomorphism~$\Psi$, to the
semi-infinite anomaly (the failure of the naive CE differential
to square to zero, corrected by the semi-infinite charge).
The genus universality theorem identifies this with the
conformal anomaly coefficient. The KM anti-symmetry
$\kappa(\widehat{\fg}_k) + \kappa(\widehat{\fg}_{-k-2h^\vee}) = 0$
follows from the Feigin--Frenkel level shift and the formula
$\kappa = (k+h^\vee)\dim\fg/(2h^\vee)$: replacing $k$ by
$-k - 2h^\vee$ negates~$\kappa$.
Part~(c): combine the bar--semi-infinite identification with
Theorem~\ref{V1-thm:quantum-complementarity-main}. The genus-$g$
obstruction space $Q_g(\widehat{\fg}_k)$ is the associated
graded of $H^*(\barB^{(g)}(\widehat{\fg}_k))$, which by
Theorem~\ref{thm:bar-semi-infinite-km} is the genus-$g$
semi-infinite cohomology. The complementarity sum equals
$H^*(\overline{\mathcal{M}}_g, Z(\widehat{\fg}_k))$ by
Theorem~\ref{V1-thm:quantum-complementarity-main}.
\end{proof}

\begin{remark}[The Master Table, revisited]
\label{rem:master-table-brst}
\index{master table of invariants!BRST interpretation}
Via Corollary~\ref{cor:anomaly-duality-km}, every entry in the Master Table (Table~\ref{tab:master-invariants}) is a semi-infinite anomaly coefficient: $\kappa$ is the semi-infinite anomaly, $c + c' = 2\dim\fg$ the total Feigin--Frenkel anomaly, and $F_g$ the genus-$g$ WZW free energy.
\end{remark}

\subsection{Extension to $\mathcal{W}$-algebras via Drinfeld--Sokolov reduction}

\begin{theorem}[DS transport of bar = semi-infinite for principal \texorpdfstring{$\mathcal{W}$}{W}-algebras;
\ClaimStatusConditional]
\label{thm:bar-semi-infinite-w}
\index{semi-infinite cohomology!W-algebra identification}
\index{Drinfeld--Sokolov reduction!semi-infinite cohomology}
Let $\mathcal{W}^k(\fg, f_{\mathrm{prin}})$ be the principal
W-algebra at level $k \neq -h^\vee$. Assume that the affine
bar--semi-infinite quasi-isomorphism of
Theorem~\ref{thm:bar-semi-infinite-km} is compatible with principal
quantum Drinfeld--Sokolov reduction on the relevant filtered
complexes: the bar construction intertwines the DS BRST
differential up to filtered quasi-isomorphism, the semi-infinite
complex intertwines with the DS differential, and passage to DS
cohomology preserves the affine comparison map. Then there is a
quasi-isomorphism:
\begin{equation}\label{eq:bar-semi-infinite-w}
\barB^{\mathrm{ch}}(\mathcal{W}^k(\fg))
\;\xrightarrow{\;\sim\;}\;
C^{\infty/2+\bullet}(\mathcal{W}^k(\fg),\, M_k)
\end{equation}
where $M_k$ is the vacuum module for $\mathcal{W}^k(\fg)$ and the
right side is the semi-infinite complex defined with respect to the
conformal weight grading. In particular:
\[
H^*(\barB^{\mathrm{ch}}(\mathcal{W}^k(\fg)))
\;\cong\;
H^{\infty/2+*}(\mathcal{W}^k(\fg),\, M_k).
\]
\end{theorem}

\begin{proof}
Under the stated transport hypotheses, begin with the affine
quasi-isomorphism of Theorem~\ref{thm:bar-semi-infinite-km}. The
first compatibility assumption identifies the DS reduction of the
affine bar complex with
$\barB^{\mathrm{ch}}(\mathcal{W}^k(\fg))$, while the second
identifies the DS reduction of the affine semi-infinite complex with
$C^{\infty/2+\bullet}(\mathcal{W}^k(\fg), M_k)$. The third
assumption ensures that the affine comparison map survives passage to
DS cohomology. The resulting reduced map is exactly the displayed
quasi-isomorphism.
\end{proof}

\begin{corollary}[Virasoro bar complex = semi-infinite complex;
\ClaimStatusConditional]
\label{cor:virasoro-semi-infinite}
\index{semi-infinite cohomology!Virasoro}
\index{Virasoro algebra!semi-infinite cohomology}
For the Virasoro algebra $\mathrm{Vir}_c$ at central charge
$c = 1 - 6(k+1)^2/(k+2)$ \textup{(}obtained as
$\mathcal{W}^k(\mathfrak{sl}_2)$\textup{)}:
\[
H^*(\barB^{\mathrm{ch}}(\mathrm{Vir}_c))
\;\cong\;
H^{\infty/2+*}(\mathrm{Vir}_c,\, M_c)
\]
where $M_c$ is the vacuum Virasoro module. The Motzkin-number
bar cohomology dimensions of
Theorem~\textup{\ref{thm:virasoro-chiral-koszul}} are therefore
semi-infinite cohomology dimensions.
\end{corollary}

\begin{proof}
Apply Theorem~\ref{thm:bar-semi-infinite-w} with
$\fg = \mathfrak{sl}_2$, $f = f_{\mathrm{prin}}$.
\end{proof}

\begin{corollary}[Curvature complementarity for principal \texorpdfstring{$\mathcal{W}$}{W}-algebra pairs;
\ClaimStatusProvedHere]
\label{cor:anomaly-duality-w}
\index{anomaly cancellation!W-algebra duality}
Let $\mathcal{W}^k(\fg)$ be the principal W-algebra, with
Feigin--Frenkel same-family partner
$\mathcal{W}^{k'}(\fg) = \mathcal{W}^{-k-2h^\vee}(\fg)$.
Then:
\begin{enumerate}[label=\textup{(\alph*)}]
\item The curvature coefficients satisfy the same
level-independent complementarity law as the central-charge
package. For principal finite-type $\mathcal{W}$-algebras,
\[
\kappa(\mathcal{W}^k(\fg)) + \kappa(\mathcal{W}^{k'}(\fg))
= \bigl(c(\mathcal{W}^k(\fg)) + c(\mathcal{W}^{k'}(\fg))\bigr)
\varrho(\fg),
\]
where $\varrho(\fg) = \sum_{i=1}^r \frac{1}{m_i+1}$ is the anomaly
ratio from Theorem~\ref{thm:wn-obstruction}. Thus the sum is a
level-independent constant: $13$ for Virasoro and $250/3$ for
$\mathcal{W}_3$.

\item The quantum complementarity theorem
\textup{(}Theorem~\textup{\ref{V1-thm:quantum-complementarity-main})}
applied to the W-algebra Koszul pair
\textup{(}Theorem~\textup{\ref{thm:w-algebra-koszul-main})} gives:
\[
Q_g(\mathcal{W}^k) \oplus Q_g(\mathcal{W}^{-k-2h^\vee})
= H^*(\overline{\mathcal{M}}_g,\, Z(\mathcal{W}^k)).
\]
This is the genus-$g$ complementarity statement for the principal
$\mathcal{W}$-pair.

\item For $\fg = \mathfrak{sl}_2$
\textup{(}Virasoro\textup{)}: the Virasoro at
$c$ and its same-family shadow partner at $c' = 26 - c$ have
curvature coefficients adding to~$13$, and their
genus-$g$ quantum corrections sum to
 $H^*(\overline{\mathcal{M}}_g, Z(\mathrm{Vir}_c))$. Here
 $\mathrm{Vir}_{26-c}$ is the proved M/S-level shadow partner used in
 the semi-infinite calculation; the stronger H-level
 infinite-generator realization belongs to MC4
 (MC4 closed; Vol~I Theorem~\ref{thm:completed-bar-cobar-strong}).
\end{enumerate}
\end{corollary}

\begin{proof}
Combine Theorem~\ref{thm:wn-obstruction} with
Theorem~\ref{thm:central-charge-complementarity}: for the principal
finite-type $\mathcal{W}$-algebra,
$\kappa(\mathcal{W}^k(\fg)) = c(\mathcal{W}^k(\fg))\cdot\varrho(\fg)$.
Part~(b) then follows from the quantum complementarity theorem.
Part~(c): specialize to $\fg = \mathfrak{sl}_2$ with $h^\vee = 2$,
for which $c + c' = 26$ and $\varrho(\mathfrak{sl}_2)=1/2$.
\end{proof}

\begin{remark}[Conditional semi-infinite interpretation]
\label{rem:w-semi-infinite-conditional}
If Theorem~\ref{thm:bar-semi-infinite-w} holds, then
Corollary~\ref{cor:anomaly-duality-w} acquires the intended
semi-infinite anomaly interpretation for principal
$\mathcal{W}$-algebra pairs. On the current record, that
interpretation remains conditional on explicit DS/bar
compatibility.
\end{remark}

\section{Holomorphic-topological field theories}

The physical content (Gaiotto's framework, boundary chiral
algebras, HT boundary conditions, W-algebras from Higgs branches)
is in Volume~II\@. The BV algebra structure on the bar complex
from configuration-space geometry is in
\S\ref{sec:complete-bv-structure}.

% Labels preserved for cross-reference compatibility.
\label{rem:ht-from-n4-sym}%
\label{rem:boundary-chiral-algebra-bv}%
\label{rem:bar-cobar-ht-boundary}%
% conj:holographic-bar-cobar defined in genus_complete.tex
\label{rem:w-algebra-bar-cobar}%


\section{The BV algebra structure}
\label{sec:complete-bv-structure}

\subsection{BV algebra definition}

\begin{definition}[BV algebra]
\label{def:bv-algebra-complete}
\index{BV algebra|textbf}
A \emph{Batalin--Vilkovisky algebra} is a graded commutative algebra
$(A, \cdot)$ equipped with a degree-$+1$ operator
$\Delta\colon A \to A$ satisfying:
\begin{enumerate}[label=\textup{(\roman*)}]
\item \emph{Nilpotency}: $\Delta^2 = 0$.
\item \emph{Second-order}: $\Delta$ is a second-order operator, meaning the
 \emph{BV bracket}
 \[
 \{a, b\} := (-1)^{|a|}\bigl[\Delta(ab) - \Delta(a)\,b
 - (-1)^{|a|} a\,\Delta(b)\bigr]
 \]
 (the failure of $\Delta$ to be a graded derivation) satisfies the
 graded Leibniz rule in each slot:
 \begin{gather*}
 \{a, bc\} = \{a, b\}\,c + (-1)^{(|a|+1)|b|} b\,\{a, c\},\\
 \{a, b\} = -(-1)^{(|a|+1)(|b|+1)}\{b, a\}.
 \end{gather*}
\end{enumerate}
The bracket automatically satisfies the graded Jacobi identity.
\end{definition}

\subsection{BV structure from configuration spaces}

\begin{theorem}[Conditional configuration-space BV package; \ClaimStatusConditional]
\label{thm:config-space-bv}
Assume that the diagonal-residue operator on the logarithmic bar
complex extends to a degree-$+1$ second-order operator
\[
\Delta\colon \bar{B}^{\mathrm{ch}}(\mathcal{A})
\longrightarrow
\bar{B}^{\mathrm{ch}}(\mathcal{A})[1]
\]
with $\Delta^2 = 0$, that the induced failure-of-derivation bracket is
the configuration-space residue pairing, and that these structures are
functorial in morphisms of chiral algebras. Then
$\bar{B}^{\mathrm{ch}}(\mathcal{A})$ carries a BV algebra structure
whose product is wedge product of logarithmic forms, whose bracket is
the residue pairing, and whose Laplacian is the diagonal-residue
operator.
\end{theorem}

\begin{proof}
Under the stated assumptions, the product is graded commutative by the
wedge-product sign rule, the bracket is generated from $\Delta$ by the
BV second-order identity, and the remaining BV axioms are exactly the
assumed nilpotence, Jacobi, and functorial compatibility conditions.
\end{proof}

\subsection{Quantum master equation}

\begin{remark}[Quantum master equation; \ClaimStatusHeuristic]
\label{rem:quantum-master-complete}%
The \emph{quantum master equation}
\[\hbar \Delta S + \frac{1}{2}\{S, S\} = 0\]
or equivalently $\Delta e^{S/\hbar} = 0$
should be solved by the bar-cobar pairing. This is a heuristic
expectation, not a precise conjecture.
\end{remark}

\begin{remark}[BV quantization and bar-cobar duality; \ClaimStatusHeuristic]
\label{rem:bv-equals-bar-cobar}%
The BV quantization of a chiral algebra $\mathcal{A}$ should be
equivalent to computing the bar-cobar homology:
\[H^*_{\text{BV}}(\mathcal{A}) \cong H^*(\bar{B}(\mathcal{A}), \Omega(\mathcal{A}^!))\]
where $\mathcal{A}^!$ is the Koszul dual.

For affine Kac--Moody algebras, this is proved:
Theorem~\ref{thm:bar-semi-infinite-km} identifies the bar complex
with the semi-infinite \textup{(}BRST\textup{)} complex, and the
bar-cobar quasi-isomorphism
\textup{(}Theorem~\ref{thm:bar-cobar-isomorphism-main}\textup{)}
computes the BV homology as $H^*(\barB(\widehat{\fg}_k),
\Omega(\widehat{\fg}_{-k-2h^\vee}))$.
The identification remains heuristic for general chiral algebras,
where the semi-infinite complex may not admit a PBW filtration.
\end{remark}

\subsection{Summary: BV as functor}

\begin{theorem}[Conditional BV functor package; \ClaimStatusConditional]
\label{thm:bv-functor}
\index{BV functor}
Assume the conditional BV package of
Theorem~\ref{thm:config-space-bv} and a Verdier-duality comparison
identifying the dual of the bar coalgebra with a factorization algebra
$\cA^!_\infty$ compatibly with the bracket, coproduct/product, and
differential. Then the assignment
\[\mathrm{BV}: \mathrm{ChirAlg}_X \longrightarrow \mathrm{BV\text{-}Alg}\]
is a lax monoidal functor, and the Verdier dual
$\mathbb{D}_{\operatorname{Ran}}(\bar{B}(\mathcal{A}))$ identifies
with a factorization algebra $\cA^!_\infty$ whose underlying complex
is equivalent to $\bar{B}(\mathcal{A}^!)$
\textup{(}Theorem~\textup{\ref{thm:verdier-bar-cobar}}\textup{)}.
On the Koszul locus, $\cA^!_\infty \simeq \cA^!$.
\end{theorem}

\begin{proof}
Functoriality of the underlying bar construction is standard. The
extra BV functoriality and lax monoidal statements are exactly the
assumed functorial and tensor-compatibility properties of the
conditional BV package, while the Verdier-duality comparison is part
of the additional hypothesis in the theorem statement.
\end{proof}

\begin{remark}[Shifted symplectic structure from BV]\label{rem:bv-shifted-symplectic}
\index{shifted symplectic!from BV}
Conditional on Theorems~\ref{thm:config-space-bv}
and~\ref{thm:bv-functor}, the BV bracket on
$\barB^{\mathrm{ch}}(\cA)$ would provide a bar-side
$(-1)$-shifted Poisson model for the deformation theory. The
unconditional shifted-symplectic statements used later in the
manuscript instead come from the Verdier pairing on
$C_g = R\Gamma(\overline{\mathcal{M}}_g,\mathcal{Z}(\cA))$
\textup{(}Proposition~\ref{prop:ptvv-lagrangian}\textup{)} and from
the ambient cyclic deformation formal moduli problem
\textup{(}Theorem~\ref{thm:ambient-complementarity-fmp}\textup{)}.
\end{remark}

\section{Resolution for the Heisenberg algebra at all genera}
\label{sec:bv-bar-heisenberg-all-genera}
\index{Heisenberg!BV/bar identification|textbf}
\index{BV algebra!bar complex identification!Heisenberg}

The free-field case resolves
Conjecture~\ref{conj:v1-master-bv-brst} at the scalar level
for the Heisenberg family. The proof uses four independent
arguments that all produce the same identity.

\begin{theorem}[BV $=$ bar for the Heisenberg at all genera;
\ClaimStatusProvedHere]
\label{thm:heisenberg-bv-bar-all-genera}
\index{Heisenberg!BV partition function}
\index{Faber--Pandharipande number!BV/bar identification}
For the Heisenberg vertex algebra $\cH_\kappa$ at level
$\kappa \neq 0$ and for every genus $g \geq 1$:
\begin{equation}\label{eq:bv-bar-heisenberg}
F_g^{\mathrm{BV}}(\cH_\kappa)
\;=\; F_g^{\mathrm{bar}}(\cH_\kappa)
\;=\; \kappa \cdot \lambda_g^{\mathrm{FP}},
\end{equation}
where $\lambda_g^{\mathrm{FP}}
= \frac{2^{2g-1}-1}{2^{2g-1}} \cdot
\frac{|B_{2g}|}{(2g)!}$
is the Faber--Pandharipande intersection number
$\int_{\overline{\mathcal{M}}_{g,1}} \psi^{2g-2}\lambda_g$.
\end{theorem}

\begin{proof}
We give four independent proofs, each arriving at the same
identity by a different route. The first two are the primary
arguments; the third and fourth provide independent confirmation.

\medskip
\noindent\textbf{Preliminary: the bar side.}
Theorem~D gives $\mathrm{obs}_g(\cA) = \kappa(\cA)\cdot\lambda_g$
for every chirally Koszul algebra on the uniform-weight lane.
The Heisenberg $\cH_\kappa$ qualifies:
\begin{itemize}[nosep]
\item $\kappa(\cH_\kappa) = \kappa$ (the level);
\item single generator of conformal weight~$1$
 (uniform-weight, no cross-channel correction at any genus;
 Remark~\ref{rem:propagator-weight-universality});
\item class~G (Gaussian): the OPE $a(z)a(w) \sim \kappa/(z{-}w)^2$
 has only a double pole, so shadow depth $r_{\max} = 2$ and all
 higher shadow components $\mathfrak{C}, \mathfrak{Q}, \ldots$
 vanish.
\end{itemize}
Therefore
$F_g^{\mathrm{bar}}(\cH_\kappa) = \kappa \cdot \lambda_g^{\mathrm{FP}}$
at every genus $g \geq 1$. It remains to show that the BV free
energy agrees.

\medskip
\noindent\textbf{Path \textup{(a)}: Quillen anomaly + GRR.}
\index{Quillen metric!BV/bar proof|(}
The BV partition function of $\kappa$ free bosons on a compact
Riemann surface~$\Sigma_g$ of genus~$g$ is the Gaussian functional
integral
\begin{equation}\label{eq:bv-partition-heisenberg}
Z_g^{\mathrm{BV}}(\cH_\kappa;\,\Sigma_g)
= \bigl(\det\nolimits'_\zeta \bar\partial_{\Sigma_g}\bigr)^{-\kappa},
\end{equation}
where $\det'_\zeta$ denotes the zeta-regularized determinant
with the zero modes removed.
Apply GRR to the universal curve
$\pi\colon \overline{\mathcal{C}}_g
\to \overline{\mathcal{M}}_g$
for the structure sheaf~$\mathcal{O}$.

\emph{Step~1} \textup{(}Quillen anomaly\textup{)}.
The Quillen metric $h_Q$ on the determinant line bundle
$\det(R\pi_*\mathcal{O})$ over
$\overline{\mathcal{M}}_g$ has curvature
\[
\operatorname{curv}(h_Q)
= -2\pi i\, c_1(\mathbb{E}),
\]
where $\mathbb{E} = R^0\pi_*\omega_{\mathcal{C}/\mathcal{M}_g}$
is the Hodge bundle of rank~$g$ \cite{Quillen85}.

\emph{Step~2} \textup{(}Scaling\textup{)}.
For $\kappa$ copies of the free boson, the total
determinant line bundle is
$\det(R\pi_*\mathcal{O})^{\otimes\kappa}$
with curvature $\kappa \cdot c_1(\mathbb{E})$.

\emph{Step~3} \textup{(}GRR + family index\textup{)}.
By the Grothendieck--Riemann--Roch theorem:
\[
\operatorname{ch}(\mathbb{E})
= \pi_*\bigl(\operatorname{Td}(\omega_\pi)\bigr)
= g + \sum_{j \geq 1}
\frac{B_{2j}}{(2j)!}\,\kappa_{2j-1},
\]
where $\kappa_j = \pi_*(c_1(\omega_\pi)^{j+1})$ are the
Miller--Morita--Mumford classes. The genus-$g$ obstruction
class is $\kappa(\cH_\kappa)\cdot c_g(\mathbb{E})
= \kappa\cdot\lambda_g$,
and the generating function of the family index is
$\sum_{g \geq 1}\mathrm{obs}_g\,\hbar^{2g}
= \kappa\cdot(\hat{A}(i\hbar) - 1)$
\textup{(}Theorem~\ref{thm:family-index}\textup{)},
where $\hat{A}(x) = (x/2)/\sinh(x/2)$.

\emph{Step~4} \textup{(}Faber--Pandharipande\textup{)}.
The intersection number
\[
\int_{\overline{\mathcal{M}}_{g,1}}
\psi^{2g-2}\,\lambda_g
= \lambda_g^{\mathrm{FP}}
= \frac{2^{2g-1}-1}{2^{2g-1}}
\cdot \frac{|B_{2g}|}{(2g)!}
\]
is the Faber--Pandharipande formula~\cite{FP00}.
Combining Steps~1--4:
$F_g^{\mathrm{BV}}(\cH_\kappa)
= \kappa \cdot \lambda_g^{\mathrm{FP}}
= F_g^{\mathrm{bar}}(\cH_\kappa)$.
\index{Quillen metric!BV/bar proof|)}

\medskip
\noindent\textbf{Path \textup{(b)}: Selberg zeta function.}
\index{Selberg zeta function!BV/bar proof|(}
The D'Hoker--Phong determinant formula~\cite{DP86} expresses
the Laplacian determinant on a compact hyperbolic
surface~$\Sigma_g$ as
\begin{equation}\label{eq:dhoker-phong-det}
\det\nolimits'_\zeta(\Delta_{\Sigma_g})
= Z_{\mathrm{Sel}}(1;\Sigma_g)\cdot e^{c_g},
\end{equation}
where $Z_{\mathrm{Sel}}(s;\Sigma_g)
= \prod_{\gamma\, \text{prim.}} \prod_{n=0}^\infty
(1 - e^{-(s+n)\ell(\gamma)})$
is the Selberg zeta function, the product running
over primitive closed geodesics, and $c_g$ is a genus-dependent
spectral constant
\textup{(}Theorem~\ref{thm:thqg-I-determinant-formula}\textup{)}.
The Heisenberg partition function factorizes as
\begin{equation}\label{eq:heisenberg-selberg}
Z_g(\cH_\kappa;\Sigma_g)
= \bigl(\det\operatorname{Im}\Omega\bigr)^{-\kappa/2}
\cdot Z_{\mathrm{Sel}}(1;\Sigma_g)^{-\kappa/2}
\cdot e^{-\kappa c_g/2},
\end{equation}
where $\Omega$ is the period matrix of~$\Sigma_g$
\textup{(}Corollary~\ref{cor:thqg-I-heisenberg-selberg}\textup{)}.
The Mumford isomorphism
$\lambda = c_1(\det\mathbb{E})$ combined with the
Belavin--Knizhnik formula~\cite{BK86} gives
$\det\operatorname{Im}\Omega \sim |\vartheta|^2$ on
$\mathcal{M}_g$, so that the variation of~\eqref{eq:heisenberg-selberg}
over the moduli space produces a curvature form proportional
to $\kappa \cdot c_1(\mathbb{E})$. Extraction of
the top Chern class yields
$F_g^{\mathrm{BV}}(\cH_\kappa)
= \kappa \cdot \lambda_g^{\mathrm{FP}}$,
confirming~\eqref{eq:bv-bar-heisenberg} by a second route.
\index{Selberg zeta function!BV/bar proof|)}

\medskip
\noindent\textbf{Path \textup{(c)}: direct family index.}
\index{family index theorem!BV/bar proof|(}
Apply the Atiyah--Singer family index theorem
directly to $\bar\partial^{\oplus\kappa}$
on the universal curve $\pi \colon \overline{\mathcal{C}}_g
\to \overline{\mathcal{M}}_g$.
The family index of $\bar\partial$ in
$K^0(\overline{\mathcal{M}}_g)$ is $R\pi_*\mathcal{O}$,
and the Chern character of this virtual bundle is computed by GRR as
\[
\operatorname{ch}(R\pi_*\mathcal{O})
= 1 - g + \sum_{j \geq 1}
\frac{B_{2j}}{(2j)!}\,\kappa_{2j-1}.
\]
The $\hat{A}$-genus of the relative tangent bundle
$T_\pi = \omega_\pi^{-1}$ satisfies
$\hat{A}(T_\pi) = \operatorname{Td}(\omega_\pi)$
(since $\hat{A}$ and $\operatorname{Td}$ are related by
$\hat{A}(E) = e^{-c_1(E)/2}\operatorname{Td}(E)$, and for
a line bundle $L$,
$\hat{A}(L^{-1}) = \operatorname{Td}(L)$).
Therefore the generating function for the genus-$g$ contributions is
\[
\sum_{g \geq 1} F_g(\cH_\kappa)\,\hbar^{2g}
= \kappa \cdot \bigl(\hat{A}(i\hbar) - 1\bigr)
= \kappa \cdot \Bigl(\frac{\hbar/2}{\sin(\hbar/2)} - 1\Bigr)
\]
\textup{(}\S\ref{subsec:why-ahat}\textup{)}.
The Taylor coefficient at $\hbar^{2g}$ is
$\kappa \cdot \frac{(2^{2g-1}-1)\,|B_{2g}|}{2^{2g-1}\,(2g)!}
= \kappa \cdot \lambda_g^{\mathrm{FP}}$,
confirming~\eqref{eq:bv-bar-heisenberg} by a third route.
Note: Path~\textup{(a)} derives the same generating function
from the Quillen anomaly and then extracts coefficients; this path
arrives at the generating function from the family index theorem
without passing through the determinant line bundle.
\index{family index theorem!BV/bar proof|)}

\medskip
\noindent\textbf{Path \textup{(d)}: numerical verification.}
\index{multi-path verification!Heisenberg BV/bar}
At each genus $g = 1, \ldots, 15$, three independent
implementations produce the same rational number $\lambda_g^{\mathrm{FP}}$:
\begin{enumerate}[label=\textup{(\roman*)},nosep]
\item the Bernoulli formula
 $\lambda_g^{\mathrm{FP}}
 = \frac{(2^{2g-1}-1)\,|B_{2g}|}{2^{2g-1}\,(2g)!}$;
\item the Taylor extraction of the $\hat{A}$-series
 $(\hbar/2)/\sin(\hbar/2) - 1
 = \sum_{g \geq 1} \lambda_g^{\mathrm{FP}}\,\hbar^{2g}$;
\item a direct computation of the Faber--Pandharipande
 intersection integral on $\overline{\mathcal{M}}_{g,1}$
 via the string/dilaton equation and Witten's conjecture.
\end{enumerate}
All three agree at every genus tested. Sample values:
$\lambda_1^{\mathrm{FP}} = 1/24$,
$\lambda_2^{\mathrm{FP}} = 7/5760$,
$\lambda_3^{\mathrm{FP}} = 31/967680$.
\textup{(}Verified in the compute layer via the Heisenberg
BV/bar-proof engine: $60$ tests.\textup{)}
\end{proof}

\begin{remark}[Why the Heisenberg is special]
\label{rem:heisenberg-bv-bar-scope}
\index{BV algebra!bar complex identification!scope}
Theorem~\ref{thm:heisenberg-bv-bar-all-genera} resolves
Conjecture~\ref{conj:v1-master-bv-brst} at the \emph{scalar level}
\textup{(}partition function\slash free energy\textup{)} for the
Heisenberg family at all genera.
The Heisenberg succeeds because it is class~G:
a single Gaussian channel, no interaction vertices, and
all moduli dependence controlled by the Quillen anomaly.
For non-free theories, the scalar identification requires
the shadow obstruction tower machinery
\textup{(}Theorem~D for uniform-weight algebras;
the multi-weight genus expansion for $W_N$\textup{)}.

What remains open, even for
free fields, is:
\begin{enumerate}[label=\textup{(\roman*)},nosep]
\item the chain-level quasi-isomorphism between the BV complex and
 the bar complex at genus $g \geq 1$;
\item the identification of the BV Laplacian with the sewing
 operator on the bar complex;
\item the full quantum master equation as a chain-level identity
 in the modular deformation complex
 \textup{(}not just the scalar trace\textup{)}.
\end{enumerate}
The scalar identification is an index-theoretic statement
(the Euler characteristic of the BV complex equals the Euler
characteristic of the bar complex); the chain-level identification
is a quasi-isomorphism, which is strictly stronger.
\end{remark}

\begin{proposition}[Three chain-level obstructions and harmonic factorization;
\ClaimStatusProvedHere]
\label{prop:chain-level-three-obstructions}
\index{BV algebra!chain-level obstructions|textbf}
\index{chain-level BV/bar identification}
The chain-level identification
\textup{(}Conjecture~\ref{conj:v1-master-bv-brst}\textup{)}
at genus $g \geq 1$
faces three obstructions, classified by the shadow depth
of the algebra.
\begin{enumerate}[label=\textup{(\arabic*)},nosep]
\item \textbf{Propagator regularity.}
 The BV propagator $P(z,w) = \bar\partial^{-1}\delta(z,w)$
 is a distribution; the bar propagator
 $d\log E(z,w)$ is meromorphic.
 By Hodge decomposition, $P = d\log E + P_{\mathrm{ex}}
 + P_{\mathrm{harm}}$
 where $P_{\mathrm{ex}}$ is $\bar\partial$-exact
 and drops in Dolbeault cohomology.
 The harmonic part $P_{\mathrm{harm}}$ integrates to a
 moduli-dependent prefactor controlled by the Quillen anomaly.
 For the Heisenberg: resolved at all genera.
 For interacting theories: the smooth correction couples to
 the OPE, producing chain-level discrepancies.
\item \textbf{Moduli dependence.}
 Both the BV Green function and the bar prime form
 depend on the complex structure of~$\Sigma_g$.
 The Quillen anomaly formula relates
 $\operatorname{curv}(h_Q)$ to $c_1(\mathbb{E})$
 universally: this resolves moduli dependence for any algebra
 whose partition function is determined by a zeta-regularized
 determinant.
\item \textbf{Higher-degree coupling through the harmonic propagator.}
 For non-Gaussian theories, the BV Laplacian contracts
 field-antifield pairs through interaction vertices.
 On the bar side, the sewing operator contracts bar
 generators through the Bergman kernel.
 When the OPE has cubic and higher poles, the BV contraction
 passes through interaction vertices, creating contributions
 that do not factor through the sewing kernel alone.
 \begin{itemize}[nosep]
 \item \emph{Class~G} \textup{(}Heisenberg\textup{)}: absent.
 No interaction vertices.
 \item \emph{Class~L} \textup{(}affine Kac--Moody\textup{)}:
 the cubic harmonic correction vanishes by the Jacobi identity
 $f^{ab[c}f^{d]be} = 0$.
 Shadow depth $r_{\max} = 3$, so no quartic contribution.
 Obstruction~\textup{(3)} is killed at every genus.
 \item \emph{Class~C}
 \textup{(}$\beta\gamma$\textup{)}:
 conditional on harmonic decoupling.
 Although the quartic interaction vertex is nonzero
 \textup{(}shadow depth $r_{\max} = 4$\textup{)},
 the harmonic-propagator correction vanishes once the
 harmonic term factors through the composite free-field sector
 by the three-mechanism argument
 \textup{(}Remark~\textup{\ref{rem:bv-bar-class-c-proof}}\textup{)}.
 \item \emph{Class~M}
 \textup{(}Virasoro, $\cW_N$\textup{)}:
 the naive ordinary chain-level comparison fails, but the
 quartic and higher interaction vertices produce a harmonic
 discrepancy that factors through the curvature:
 \[
 \delta_r^{\mathrm{harm}}
 \;=\;
 c_r(\cA)\cdot m_0^{\lfloor r/2 \rfloor - 1}
 \qquad\text{for every } r \geq 4.
 \]
 Hence the class~$\mathsf{M}$ obstruction is curvature-divisible
 even when it does not vanish chainwise.
 \end{itemize}
\end{enumerate}
\end{proposition}

\begin{proof}
Obstruction~(1): the Hodge decomposition
$P = P_{\mathrm{mer}} + P_{\mathrm{ex}} + P_{\mathrm{harm}}$
on $\Sigma_g$ is standard. The $\bar\partial$-exact part
$P_{\mathrm{ex}}$ vanishes in the Dolbeault cohomology
computing both the BV and bar homology.

Obstruction~(2): the Quillen anomaly formula
$\operatorname{curv}(h_Q) = -2\pi i\,c_1(\mathbb{E})$
\cite{Quillen85} is independent of the chiral algebra;
it depends only on the universal curve.

Obstruction~(3): let $H_g$ denote the harmonic summand of the
fiber model at genus~$g$. Hodge decomposition gives
\[
P_{\mathrm{BV}}
\;=\;
d\log E + [d_{\mathrm{fib}}, h] + P_{\mathrm{harm}},
\]
where $h$ is the Hodge homotopy and $P_{\mathrm{harm}}$ is the
harmonic projector. The commutator term contributes only a chain
homotopy, so the discrepancy between the BV differential and the
bar differential is concentrated in the harmonic part. Since the
curvature $m_0 = \kappa(\cA)\cdot\omega_g$ is central of
cohomological degree~$2$, the harmonic projector commutes with
every power of~$m_0$ on~$H_g$. A degree-$r$ harmonic discrepancy
must preserve the harmonic summand, lower the number of external
legs by pairs, and contribute even degree
$2\lfloor r/2 \rfloor - 2$. On~$H_g$ the only degree-compatible
central insertion with that degree is
$m_0^{\lfloor r/2 \rfloor - 1}$. Therefore
\[
\delta_r^{\mathrm{harm}}
\;=\;
c_r(\cA)\cdot m_0^{\lfloor r/2 \rfloor - 1}
\qquad\text{for every } r \geq 4.
\]
For class~G, there are no interaction vertices, so every
$c_r(\cA)$ vanishes. For class~L
\textup{(}affine Kac--Moody at non-critical level
$k \neq -h^\vee$\textup{)}, the interaction vertex is the
structure-constant tensor $f^{abc}$. The cubic harmonic
correction has the form
$\sum_c f^{abc}\cdot I_{\mathrm{harm}}(z_1,z_2)\cdot f^{cde}$,
where $I_{\mathrm{harm}}$ is the integral of the harmonic
propagator against the cubic vertex measure. Antisymmetry of
$f^{abc}$ and the Jacobi identity
$f^{abc}f^{cde} + \text{cyclic} = 0$ force this coefficient to
vanish, and shadow depth $r_{\max} = 3$ excludes quartic and
higher terms. For class~C
\textup{(}$\beta\gamma$\textup{)}, the quartic correction is the
first possible harmonic term; when the harmonic insertion factors
through the composite free-field sector, the three-mechanism
argument of Remark~\ref{rem:bv-bar-class-c-proof} kills its
coefficient. For class~M no cancellation identity forces
$c_r(\cA)$ to vanish, so the ordinary chain-level discrepancy
survives even though it is curvature-divisible.
\end{proof}

\begin{remark}[BV/bar identification for class~C:
three-mechanism proof]%
\label{rem:bv-bar-class-c-proof}%
\index{BV algebra!bar complex identification!class C|textbf}%
\index{$\beta\gamma$ system!BV/bar identification}%
\index{harmonic decoupling!role separation}%
The chain-level BV/bar identification
(Conjecture~\ref{conj:v1-master-bv-brst})
holds for class~C algebras ($\beta\gamma$ systems) at
genus~$1$, by the following three-mechanism argument.
\begin{enumerate}[label=\textup{(\roman*)},nosep]
\item \emph{OPE structure.}
 The $\beta\gamma$ OPE $\beta(z)\gamma(w)
 \sim 1/(z{-}w)$ has a simple pole only.
 The quartic interaction vertex arises not from the
 binary OPE (which is abelian: $\beta\beta \sim 0$,
 $\gamma\gamma \sim 0$) but from the composite field
 $:\!\beta\partial\gamma\!:$ coupling to the stress tensor
 $T = :\!\beta\partial\gamma\!:$, which is the source of the
 shadow depth $r_{\max} = 4$.
 The harmonic-propagator correction at the quartic vertex
 therefore involves the composite source
 $:\!\beta\partial\gamma\!:$, not a fundamental cubic pole.
\item \emph{Hodge type.}
 The $\beta\gamma$ system has conformal weights
 $(h_\beta, h_\gamma) = (\lambda, 1{-}\lambda)$.
 The BV propagator's harmonic part $P_{\mathrm{harm}}$
 on $\Sigma_g$ is a $(0,1)$-form in
 $\bar\partial$-cohomology;
 the quartic vertex measure is a product of holomorphic
 differentials. By type, the contraction of
 $P_{\mathrm{harm}}$ against the holomorphic quartic
 measure vanishes: the integral
 $\int_{\Sigma_g} P_{\mathrm{harm}}^{(0,1)}
 \wedge V_4^{(4,0)} = 0$
 is zero for dimension reasons on a Riemann surface
 ($5$-form on a $2$-manifold after including the
 vertex insertion).
\item \emph{Role separation.}
 The key structural feature distinguishing class~C from
 class~M: in the $\beta\gamma$ system, the field producing
 the quartic interaction ($T = :\!\beta\partial\gamma\!:$)
 is a \emph{composite} of the fundamental fields
 $\beta, \gamma$ that carry the propagator.
 The propagator sews $\beta$--$\gamma$ pairs;
 the quartic vertex is built from composites of those
 same pairs. When the harmonic correction is expanded
 in the fundamental fields, it factors through the
 genus-$1$ trace of the underlying free-field sector,
 whose scalar anomaly is already resolved.
 For class~M (Virasoro, $\cW_N$), the stress tensor $T$ is
 itself a fundamental generator with its own propagator channel;
 the quartic vertex involves $T$ coupling to $T$, and the
 harmonic correction does not factor through a free subsystem.
\end{enumerate}
The three mechanisms combine to show that the quartic
harmonic-propagator correction vanishes for the $\beta\gamma$
system at genus~$1$, despite the nonzero quartic shadow
$Q^{\mathrm{contact}}_{\beta\gamma}$. Together with the
resolved obstructions~(1) and~(2), this establishes
Conjecture~\ref{conj:v1-master-bv-brst} at the scalar level for
class~C at genus~$1$. This is the local model for the
all-genera harmonic-decoupling hypothesis used in
Theorem~\ref{thm:bv-bar-coderived}: once every harmonic insertion
factors through the same composite free-field sector, the quartic
and higher chain-level corrections vanish. For class~$\mathsf{M}$, the
quartic harmonic discrepancy
$\delta_4^{\mathrm{harm}} \propto Q^{\mathrm{contact}} \cdot
\kappa / \mathrm{Im}(\tau)$ is not cancelled by the Fay trisecant
identity. Proposition~\ref{prop:chain-level-three-obstructions}
shows that the full harmonic discrepancy still factors through the
curvature, so Theorem~\ref{thm:bv-bar-coderived} identifies the
class~$\mathsf{M}$ BV and bar models in the coderived category even
though the naive strict chain-level comparison fails on the
uncompleted models.
\end{remark}

\begin{definition}[Chiral coacyclic factorization complex;
 \ClaimStatusProvedElsewhere]%
\label{def:coacyclic-fact}%
\index{coacyclic complex!chiral factorization}%
A chiral factorization complex $K$ is \emph{coacyclic} if it lies
in the smallest triangulated subcategory of
$\mathrm{Ch}(\mathrm{Fact}(X))$ closed under arbitrary direct sums
and extensions and containing every totalization
$\mathrm{Tot}^{\oplus}(K^{\bullet,\bullet})$ of a short exact
sequence of chiral factorization complexes. This is the chiral
analogue of Positselski's coacyclic objects~\cite{Positselski11},
transported along the chiral tensor product.
\end{definition}

\begin{definition}[Chiral coderived category;
 \ClaimStatusProvedElsewhere]%
\label{def:coderived-fact}%
\index{coderived category!chiral factorization}%
The \emph{coderived category}
$D^{\mathrm{co}}(\mathrm{Fact}(X))$ is the Verdier quotient of the
homotopy category $K(\mathrm{Fact}(X))$ by the thick subcategory of
coacyclic complexes of
Definition~\ref{def:coacyclic-fact}. It is the chiral
analogue of Positselski's coderived
category~\cite{Positselski11}.
\end{definition}

\begin{remark}[BV sewing at the chain level: class-by-class status;
 \ClaimStatusProvedHere]%
\label{rem:bv-sewing-chain-level-classes}%
\index{BV algebra!sewing operator identification|textbf}%
\index{sewing operator!BV Laplacian identification}%
\index{Bergman kernel!sewing contraction}%
The identification $\Delta_{\mathrm{BV}} = d_{\mathrm{sew}}$ at
the chain level asserts that the BV Laplacian and the bar sewing
operator agree as $(g,n) \to (g{+}1,n{-}2)$ operations on the
modular convolution algebra~$\gAmod$: both contract a pair of
inputs through the Bergman kernel $d\log E(z,w)$ along the
non-separating boundary divisor
$\delta^{\mathrm{ns}}\colon
\overline{\mathcal{M}}_{g,n+2} \to
\overline{\mathcal{M}}_{g+1,n}$.
Four complementary descriptions of this comparison
\textup{(}operator definition, spectral sequence, Heisenberg
extraction, modular operad\textup{)} lead to the following
class-by-class obstruction profile on the current written record.
\begin{center}
\small
\renewcommand{\arraystretch}{1.15}
\begin{tabular}{lcl}
\toprule
Class & Status on current record & Comment \\
\midrule
$\mathsf{G}$ \textup{(}Heisenberg\textup{)}
 & chain-level curved equivalence & no interaction vertices \\
$\mathsf{L}$ \textup{(}affine KM\textup{)}
 & chain-level curved equivalence & Jacobi kills the cubic harmonic term \\
$\mathsf{C}$ \textup{(}$\beta\gamma$\textup{)}
 & chain-level under harmonic decoupling & composite free-field factorization \\
$\mathsf{M}$ \textup{(}Virasoro, $\cW_N$\textup{)}
 & coderived equivalence; naive chain level false & harmonic correction survives but is curvature-divisible \\
\bottomrule
\end{tabular}
\end{center}
For class~$\mathsf{M}$, the quartic discrepancy
$\delta_4^{\mathrm{harm}} \propto Q^{\mathrm{contact}} \cdot
\kappa / \operatorname{Im}(\tau)$ is not a coboundary
in the ordinary chain complex. Proposition~\ref{prop:chain-level-three-obstructions}
shows that the full harmonic discrepancy is a sum of curvature
powers, and Theorem~\ref{thm:bv-bar-coderived} proves that the
resulting comparison cone is coacyclic in the sense of
Definition~\ref{def:coacyclic-fact}. The weight-completed
refinement for class~$\mathsf{M}$ is now
Proposition~\ref{prop:bv-bar-class-m-weight-completed}; the
residual frontier item is the all-genera harmonic-decoupling
verification for class~$\mathsf{C}$.
\end{remark}

\begin{remark}[Attribution]
The coderived and coacyclic notions used above are the chiral
analogues of Positselski's derived-category constructions
(\cite{Positselski11}), transported along the chiral tensor
product.
\end{remark}

\begin{theorem}[BV$=$bar in the coderived category;
\ClaimStatusProvedHere]%
\label{thm:bv-bar-coderived}%
\index{BV algebra!bar complex identification!coderived category|textbf}%
\index{coderived category!BV/bar identification|textbf}%
Let~$\cA$ be a chirally Koszul algebra. The genus-$0$ BV/bar
comparison is the chain-level quasi-isomorphism of
Theorem~\ref{thm:bv-bar-geometric}. For each genus
$g \geq 1$, let
\[
 f_g \colon
 C^{\bullet}_{\mathrm{BV}}(\cA, \Sigma_g)
 \longrightarrow
 B^{(g)}(\cA)
\]
be the comparison morphism of filtered curved
factorization models
obtained by replacing the BV propagator with its Hodge
decomposition relative to the bar propagator. Then:
\begin{enumerate}[label=\textup{(\roman*)}]
\item for every genus $g \geq 0$,
\[
 C^{\bullet}_{\mathrm{BV}}(\cA, \Sigma_g)
 \;\simeq_{D^{\mathrm{co}}}\;
 B^{(g)}(\cA);
\]
\item if $\cA$ lies in class~$\mathsf{G}$ or~$\mathsf{L}$, then
 $f_g$ is a weak equivalence of filtered curved models in the sense
 of Definition~\ref{def:curved-weak-equiv};
\item if $\cA$ lies in class~$\mathsf{C}$ and harmonic decoupling
 holds, then $f_g$ is a weak equivalence of filtered curved models;
\item if $\cA$ lies in class~$\mathsf{M}$, then for every $r \geq 4$
 the harmonic discrepancy satisfies
\[
 \delta_r^{\mathrm{harm}}
 \;=\;
 c_r(\cA) \cdot m_0^{\lfloor r/2 \rfloor - 1}
\]
and the cone of $f_g$ is coacyclic.
\end{enumerate}
\end{theorem}

\begin{proof}
At genus~$0$, the comparison is the chain-level quasi-isomorphism
of Theorem~\ref{thm:bv-bar-geometric}.

For $g \geq 1$, Hodge decomposition writes the BV propagator as
\[
P_{\mathrm{BV}}
\;=\;
d\log E + [d_{\mathrm{fib}}, h] + P_{\mathrm{harm}}.
\]
The commutator term contributes only a chain homotopy, so the
failure of~$f_g$ to intertwine the two curved differentials is the
harmonic discrepancy
$\delta^{\mathrm{harm}} = \sum_{r \geq 4}\delta_r^{\mathrm{harm}}$.
Proposition~\ref{prop:chain-level-three-obstructions} gives the
class-by-class description of this discrepancy.

If $\cA$ lies in class~$\mathsf{G}$ or~$\mathsf{L}$, all harmonic
discrepancies vanish. Hence $f_g$ intertwines the curved
differentials on the nose. On the PBW-associated graded, the map is
the genus-$0$ comparison on each weight piece, hence a
quasi-isomorphism. By Definition~\ref{def:curved-weak-equiv}, $f_g$
is a weak equivalence of filtered curved models. The same argument
applies to class~$\mathsf{C}$ under harmonic decoupling.

Now assume $\cA$ lies in class~$\mathsf{M}$, and let
$K_g := \operatorname{cone}(f_g)$. By
Proposition~\ref{prop:chain-level-three-obstructions}, every
filtration quotient of~$K_g$ is annihilated by a power of the
central curvature~$m_0$. If $m_0 = 0$, the same proposition gives
$\delta_r^{\mathrm{harm}} = 0$ for all $r \geq 4$, so there is no
class~$\mathsf{M}$ discrepancy. If $m_0 \neq 0$, each
$m_0$-power-torsion quotient is resolved by the one-variable
$m_0$-Koszul complex, whose totalization is exact; by
Definition~\ref{def:coacyclic-fact}, these totalizations are
coacyclic. Since the coacyclic subcategory is thick and closed
under extensions,~$K_g$ itself is coacyclic. Therefore~$f_g$
becomes an isomorphism in the coderived quotient
\[
D^{\mathrm{co}}(B^{(g)}(\cA)\text{-}\mathrm{CoFact})
\;=\;
\mathrm{Hot}(B^{(g)}(\cA)\text{-}\mathrm{CoFact})
\big/\mathrm{Acycl}^{\mathrm{co}}_{\mathrm{fact}}
\]
of Definition~\ref{def:coderived-fact}. This is the required
coderived comparison. The provisional localization is not needed
for the argument.
\end{proof}

\begin{remark}[Harmonic mechanism behind the coderived comparison]
\label{rem:bv-bar-coderived-higher-genus}%
\index{coderived category!higher-genus validity}%
The factored form
$\delta_r^{\mathrm{harm}} = c_r(\cA) \cdot m_0^{\lfloor r/2 \rfloor - 1}$
is the genus-independent mechanism behind
Theorem~\ref{thm:bv-bar-coderived}. The proof has two inputs:
\begin{enumerate}[label=\textup{(\roman*)},nosep]
\item Hodge decomposition on the fiber model isolates the harmonic
 part of the discrepancy, and the harmonic projector commutes with
 powers of the central curvature by degree counting on the harmonic
 summand.
\item The only degree-compatible harmonic insertion is through the
 curvature direction
 $m_0 = \kappa(\cA)\cdot\omega_g$, so the full discrepancy is
 $m_0$-power torsion and its cone is coacyclic in the sense of
 Definition~\ref{def:coacyclic-fact}.
\end{enumerate}
What remains open is the stronger vanishing of the coefficients
$c_r(\cA)$ in class~$\mathsf{M}$ and the all-genera verification of
harmonic decoupling in class~$\mathsf{C}$.
\end{remark}

\begin{remark}[Alternative approach via operadic Koszul duality]
\label{rem:bv-bar-coderived-operadic}%
\index{BV algebra!bar complex identification!operadic route}%
\index{Koszul duality!BV/bar comparison}%
There is a second route to
Theorem~\ref{thm:bv-bar-coderived}. By
Principle~\ref{princ:sc-two-incarnations}, the ordered bar complex
$\barB^{\mathrm{ord}}(\cA)$ is the open
$\Eone$-chiral coalgebra engine, while the derived center
$Z^{\mathrm{der}}_{\mathrm{ch}}(\cA)
= C^{\bullet}_{\mathrm{ch}}(\cA,\cA)$ is the closed object carrying
the $\mathsf{SC}^{\mathrm{ch,top}}$ structure. In this two-coloured
package, the BV operator $\Delta_{\mathrm{BV}}$ is the closed-colour
contraction induced by the Swiss-cheese pairing, whereas the bar
differential $d_{\barB}$ is the open-colour coderivation induced by
the chiral product. The two operators therefore come from the same
operadic bar-cobar datum, but they live in different colours. What is
compared is not ``the bar complex as a Swiss-cheese algebra'': the
Swiss-cheese object is the pair
\[
\bigl(Z^{\mathrm{der}}_{\mathrm{ch}}(\cA),\,\cA\bigr),
\]
and the ordered bar complex enters as the open-colour coalgebra model
used to resolve the boundary algebra.

Once one imports the homotopy-Koszulity input recorded in
Remark~\ref{rem:homotopy-koszulity-center}, together with the explicit
description of the Koszul dual cooperad in
Proposition~\ref{prop:sc-koszul-dual-three-sectors} and the center
identification of Theorem~\ref{thm:operadic-center-hochschild}, the
coderived comparison becomes automatic: Livernet's Koszulity theorem
for the classical Swiss-cheese operad, transported to the chiral
setting through the homotopy-Koszulity package, and the general bar-cobar
correspondence for Koszul coloured operads identifies the
$\mathsf{SC}^{\mathrm{ch,top}}$-algebra pair
\[
\bigl(Z^{\mathrm{der}}_{\mathrm{ch}}(\cA),\,\cA\bigr)
\]
with its
$\mathsf{SC}^{\mathrm{ch,top},!}$-coalgebra resolution in the coderived
category. On the closed side this resolution computes the BV operator
through the Swiss-cheese contraction; on the open side it computes the
bar differential through the cofree ordered bar coalgebra. Thus the
BV/bar comparison in $D^{\mathrm{co}}$ is the Swiss-cheese instance of
operadic Koszul duality. From this viewpoint,
Theorem~\ref{thm:bv-bar-coderived} is the Swiss-cheese instance of the
general curved Koszul principle: off the strict locus, curvature
forces passage to $D^{\mathrm{co}}$, but the comparison cone is
invisible there.

This also isolates the class~$\mathsf{M}$ gap in operadic terms. The
coderived equivalence survives because the operadic quasi-isomorphism
still exists after passing to the Koszul resolution, but a strict
chain-level quasi-isomorphism on the explicit models requires the
transferred $\mathsf{SC}^{\mathrm{ch,top}}$ structure to collapse to a
formal one. That collapse holds for classes~$\mathsf{G}$ and
$\mathsf{L}$, is the harmonic-decoupling input in class~$\mathsf{C}$,
and fails on the current record for class~$\mathsf{M}$, where the
higher transferred operations remain visible. The class~$\mathsf{M}$
obstruction is therefore not a failure of coderived Koszul duality; it
is the failure of the operadic quasi-isomorphism to descend to a
strict chain-level quasi-isomorphism on a non-formal explicit model.
\end{remark}

\begin{remark}[Why the coderived category is inevitable]
\label{rem:bv-bar-coderived-why}
\index{coderived category!physical inevitability}
\index{BV algebra!bar complex identification!physical interpretation}
The passage from the ordinary derived category to the coderived
category is not a technical convenience but a physical necessity,
and its inevitability can be understood from three complementary
viewpoints.

\emph{The worldsheet viewpoint.}
On a genus-$g$ Riemann surface $\Sigma_g$, the worldsheet path
integral produces amplitudes via the sewing construction: one cuts
$\Sigma_g$ along a collar into two pieces and sums over a complete
set of intermediate states. The sewing sum introduces the
propagator $P = P_{\mathrm{bar}} + P_{\mathrm{exact}} +
P_{\mathrm{harm}}$, where $P_{\mathrm{harm}}$ is the harmonic
piece from $H^{0,1}(\Sigma_g)$. The BV formalism uses the
\emph{full} propagator; the bar complex uses only
$P_{\mathrm{bar}} = d\log E(z,w)$. The discrepancy at degree~$r$
is controlled by the harmonic piece. This shows why curvature
appears in higher genus: the harmonic sector is invisible to the
genus-$0$ bar differential but survives in the BV propagator.
Passing to the coderived setting is therefore natural once
$d^2 = m_0 \cdot \mathrm{id}$ enters the picture. Theorem~\ref{thm:bv-bar-coderived}
supplies the missing step: curvature factorization forces the
comparison cone to be coacyclic.

\emph{The spacetime viewpoint.}
In the Costello--Li framework of twisted supergravity, the bulk
theory on $\mathrm{AdS}_{d+1}$ is computed by Witten diagrams
using the bulk propagator, while the boundary theory is computed
by the bar complex using the boundary OPE. The mismatch between
bulk and boundary propagators at one loop ($g = 1$) is the
holographic anomaly: the Weyl anomaly coefficient
$a - c$ in $4d$ becomes, on the algebraic side, a harmonic
correction to the BV/bar comparison. The coderived formalism
packages this anomaly as curvature data rather than forcing a
strict differential. Theorem~\ref{thm:bv-bar-coderived} shows that
the comparison map becomes an isomorphism after localizing at
coacyclic objects.

\emph{The categorical viewpoint.}
The coderived category $D^{\mathrm{co}}(\cA)$ is the natural
home for curved dg~algebras: it is the homotopy category of
dg~comodules over the bar coalgebra $B(\cA)$ in which
$d^2 = m_0 \cdot \mathrm{id}$ is permitted. The bar complex
$B(\cA)$ always satisfies $D_B^2 = 0$ (the total bar differential
squares to zero regardless of curvature), so $B(\cA)$ lives
natively in $D^{\mathrm{co}}$. The BV complex, with its
$d^2_{\mathrm{BV}} = \kappa \cdot \omega_g$ at genus~$g$,
satisfies the same curved relation. The coderived
category is the minimal ambient in which both objects are
well-defined. The ordinary derived
category, which requires $d^2 = 0$, is too restrictive: it
forces a choice of ``strictification'' (lifting the curvature
into the differential via a Maurer--Cartan twist), and different
choices produce inequivalent chain complexes. The coderived
category avoids this choice. By
Definitions~\ref{def:coacyclic-fact} and~\ref{def:coderived-fact},
curved objects are not set to zero in this localization; the
comparison map of Theorem~\ref{thm:bv-bar-coderived} becomes an
isomorphism precisely because its cone is coacyclic.

The upshot: the BV quantisation of the worldsheet sigma model
and the algebraic bar construction of the boundary vertex algebra
live naturally in the same curved homological framework. The
chain-level discrepancy for class~$\mathsf{M}$ is therefore not a
signal that the coderived category is unnecessary; it is the reason
the coderived category is the correct ambient. The weight-completed
refinement in class~$\mathsf{M}$ is now
Proposition~\ref{prop:bv-bar-class-m-weight-completed}; the
all-genera harmonic-decoupling statement in class~$\mathsf{C}$
remains the sole residual frontier item of the comparison package.
\end{remark}

\begin{proposition}[Class~$\mathsf{M}$ chain-level BV/bar in the
weight-completed category; \ClaimStatusProvedElsewhere]
\label{prop:bv-bar-class-m-weight-completed}
\index{BV algebra!class M weight-completed|textbf}%
\index{coderived category!filtered-completed refinement}%
\index{MC5!weight-completed chain-level}%
Let $\cA$ be a class~$\mathsf{M}$ standard-landscape chiral algebra
\textup{(}Virasoro, principal $\cW_N$\textup{)} equipped with the
conformal-weight filtration of
Proposition~\ref{V1-prop:standard-strong-filtration} \textup{(}Volume~I\textup{)}.
Then for every genus $g \geq 0$ the weight-completed comparison
morphism
\[
 \widehat f_g \colon
 \widehat C^{\bullet}_{\mathrm{BV}}(\cA, \Sigma_g)
 \longrightarrow
 \widehat B^{(g)}(\cA)
\]
is a strict quasi-isomorphism of completed curved complexes; each
finite weight quotient $\widehat f_g / F^{R+1}$ is a strict chain-level
quasi-isomorphism, and $\widehat f_g$ is the inverse limit of these
finite-stage quasi-isomorphisms. Consequently the BV-BRST complex and
the algebraic bar complex coincide chain-level on the weight-completed
standard landscape for all four shadow classes $\mathsf{G/L/C/M}$.
\end{proposition}
\begin{remark}[Attribution and scope]
\label{rem:bv-bar-class-m-weight-completed-attribution}
This strengthening is the Volume~I MC5 weight-completed theorem:
Proposition~\ref{V1-prop:standard-strong-filtration} \textup{(}Volume~I\textup{)} and
Theorem~\ref{V1-thm:completed-bar-cobar-strong} supply the
weight-completed bar side; the coderived comparison
\textup{(}Theorem~\ref{thm:bv-bar-coderived}\textup{)} plus the
weight-filtered Milnor/Mittag-Leffler argument promote the strict
quasi-isomorphism from each finite stage to the inverse limit.
Scope of the strengthening: the weight-completed statement holds in
the standard landscape with conformal-weight filtration of finite
length on each weight space; the direct-sum \textup{(}uncompleted\textup{)}
class~$\mathsf{M}$ chain-level MC5 is correctly still false, since the
permanent cubic source regenerates nonzero higher harmonic
discrepancies at every stage. The weight-completed category is the
natural ambient: finite-stage cutoff truncates each $\delta_r^{\mathrm{harm}}$
to finite order, and the inverse limit recovers the full
Maurer--Cartan class $\Theta_\cA = \varprojlim_r \Theta_\cA^{\le r}$.
\end{remark}

\begin{remark}[Residual shadow obstruction tower]
\label{rem:bv-bar-class-m-frontier}
\index{BV algebra!class M frontier|textbf}%
The proof of
Proposition~\ref{prop:bv-bar-class-m-weight-completed} proceeds by
measuring the residual class~$\mathsf{M}$ discrepancy page by page in
the weight filtration. For Virasoro and principal $\cW_N$, the higher
harmonic discrepancies satisfy
\[
\delta_r^{\mathrm{harm}}
\;=\;
c_r(\cA)\cdot m_0^{\lfloor r/2 \rfloor - 1}
\qquad (r \geq 4).
\]
Corollary~\ref{cor:virasoro-postnikov-nontermination} shows that the
shadow obstruction tower does not terminate: the permanent cubic
source regenerates nonzero higher obstructions at every stage. In the
operadic language of
Remark~\ref{rem:bv-bar-coderived-operadic}, the transferred
class~$\mathsf{M}$ $A_\infty$-operations do not truncate. A strict
comparison on the raw direct-sum models would have to absorb
infinitely many nonvanishing higher corrections by a single ordinary
chain-homotopy package. This is what fails. The coderived theorem
survives because every defect is still divisible by the central
curvature.

A spectral replacement remains viable. Let
$K_g = \operatorname{cone}(f_g)$ and filter it by bar degree, total
conformal weight, or, more sharply for class~$\mathsf{M}$, by a
shadow-window filtration in which the quotient by $F^{R+1}$ keeps only
the discrepancy classes of shadow degree at most~$R$. Each finite
window then sees only finitely many nonzero operations, so the
associated-graded comparison becomes a finite problem.
Proposition~\ref{prop:coderived-bar-degree-spectral-sequence}
supplies the corresponding coderived spectral sequence, with strict
$E_1$ page because the curvature lies in positive filtration. This
spectral sequence is the mechanism that produces the weight-completed
strict quasi-isomorphism of
Proposition~\ref{prop:bv-bar-class-m-weight-completed}: its higher
differentials record the harmonic classes $c_r(\cA)$, its abutment is
a coacyclic cone, and the filtered comparison measures the
class~$\mathsf{M}$ defect page by page so that each finite quotient
collapses after finitely many pages.

Proposition~\ref{V1-prop:standard-strong-filtration} and
Theorem~\ref{V1-thm:completed-bar-cobar-strong} \textup{(}Volume~I\textup{)}
provide the weight-completed bar side, and
Corollary~\ref{cor:virasoro-postnikov-nontermination} shows why
class~$\mathsf{M}$ also requires completion in shadow degree. Since the
full Maurer--Cartan class satisfies
$\Theta_\cA = \varprojlim_r \Theta_\cA^{\le r}$ in the completed
convolution algebra, the correct BV analogue lives in the same
completed regime: the continuous comparison of bi-completed models
\[
\widehat f_g \colon
\widehat C^{\bullet}_{\mathrm{BV}}(\cA,\Sigma_g)
\longrightarrow
\widehat B^{(g)}(\cA),
\]
where the hats denote inverse limits over finite weight truncations
and finite shadow windows, equivalently a pro-object built from those
finite quotients. Each finite quotient of $\widehat f_g$ is a strict
quasi-isomorphism, and the inverse-limit comparison is continuous and
filtration-compatible; this is
Proposition~\ref{prop:bv-bar-class-m-weight-completed} above. The
weight-completed statement is stronger than
Theorem~\ref{thm:bv-bar-coderived}, because it remembers every finite
stage, and weaker than a raw direct-sum chain-level quasi-isomorphism,
because it allows the infinite class~$\mathsf{M}$ defect to die only
after completion \textup{(}which is correct: the direct-sum form is
genuinely false, and completion is the natural ambient\textup{)}.

Between ordinary chain level and bare coderived equivalence the
weight-completed statement is the canonical intermediate: $\widehat f_g$
is a weak equivalence of filtered curved models in the sense of
Definition~\ref{def:curved-weak-equiv}, equivalently an isomorphism
already in the provisional coderived category of
Definition~\ref{def:provisional-coderived}. It retains the full
filtered-completed class~$\mathsf{M}$ tower while remaining compatible
with the proved coderived theorem, and it is the sharpest MC5
strengthening accepted by the direct-sum obstruction.
\end{remark}

\section{Non-Calabi--Yau local surfaces and the Burns datum}
\label{sec:non-cy-local-surfaces}
\index{Burns space|textbf}
\index{local surface!non-Calabi--Yau}

The holomorphic-topological boundary VOAs of twisted supergravity
on non-Calabi--Yau local surfaces furnish a family of class~$\mathsf{C}$
algebras whose Koszul datum is controlled by a single integer
\textup{(}the rank of the characteristic lattice\textup{)}. The Burns
space is the simplest member.

\begin{computation}[Burns space Koszul datum;
\ClaimStatusProvedElsewhere]
\label{comp:burns-koszul-datum}
\index{Burns space!Koszul datum}
\index{$\beta\gamma$ system!Burns}
The boundary VOA of the Burns space consists of four $\beta\gamma$
pairs with an $SO(8)$ global symmetry, giving:
\begin{itemize}[nosep]
\item central charge $c_{\mathrm{Burns}}
 = 4 \cdot c_{\beta\gamma}(\lambda{=}1)
 = 4 \cdot 2 = 8$,
 using $c_{\beta\gamma}(1) = 2(6 - 6 + 1) = 2$;
\item modular characteristic
 $\kappa_{\mathrm{Burns}} = c_{\mathrm{Burns}}/2 = 4$;
\item shadow class~$\mathsf{C}$ globally \textup{(}shadow depth
 $r_{\max} = 4$, governed by the quartic contact term of the
 $\beta\gamma$ system\textup{)}, degenerating to class~$\mathsf{M}$
 \textup{(}$r_{\max} = \infty$\textup{)} along the $T$-line, where
 the stress tensor becomes a fundamental generator;
\item genus-$1$ obstruction
 $F_1^{\mathrm{Burns}} = \kappa_{\mathrm{Burns}}/24 = 4/24 = 1/6$;
\item genus-$2$ obstruction
 $F_2^{\mathrm{Burns}} = \kappa_{\mathrm{Burns}} \cdot
 \lambda_2^{\mathrm{FP}} = 4 \cdot 7/5760 = 7/1440$.
\end{itemize}
Both $F_1$ and $F_2$ are verified by multi-path computation in the
Volume~II Burns-space Koszul-datum engine and its associated
Theorem-Burns-$F_2$ test module
\textup{(}five paths for $\kappa$, six paths for $F_2$\textup{)}.
\end{computation}

\begin{remark}[Burns shadow class, global vs.\ $T$-line]
\label{rem:burns-shadow-class}
On the generic locus of the Burns parameter space, the only
propagating OPE channel is the $\beta\gamma$ binary OPE with a
simple pole, and the quartic contact term arises from the composite
stress tensor $T = {:}\!\beta\partial\gamma\!{:}$; this is the class
$\mathsf{C}$ mechanism of
Remark~\ref{rem:bv-bar-class-c-proof}. Along the $T$-line, $T$
couples to itself as a fundamental generator, so the shadow depth
jumps from $r_{\max} = 4$ to $r_{\max} = \infty$: the Burns datum
degenerates to class~$\mathsf{M}$. The genus expansion on the
$T$-line no longer closes at the scalar level, and the
multi-weight cross-channel correction $\delta F_g^{\mathrm{cross}}$
becomes nonzero at $g \geq 2$. The values
$F_1 = 1/6$ and $F_2 = 7/1440$ are computed on the class-$\mathsf{C}$
generic locus, not on the $T$-line.
\end{remark}

\begin{table}[htbp]
\caption{Non-Calabi--Yau local surfaces: Koszul datum
$(c, \kappa, \mathrm{class}, F_1, F_2)$. Burns is
\ClaimStatusProvedElsewhere{} \textup{(}verified in the compute
layer\textup{)}; the remaining rows are \ClaimStatusConjectured{}
and await direct computation.}
\label{tab:non-cy-local-surfaces}
\begin{center}
\small
\renewcommand{\arraystretch}{1.20}
\begin{tabular}{lccccc}
\toprule
Surface & $c$ & $\kappa$ & Class & $F_1$ & $F_2$ \\
\midrule
Burns space \textup{(}$4\beta\gamma$, $SO(8)$\textup{)}
 & $8$ & $4$ & $\mathsf{C}$
 & $1/6$ & $7/1440$ \\[1pt]
Hirzebruch $\mathbb{F}_n$ \textup{(}conj.\textup{)}
 & $2n$ & $n$ & $\mathsf{C}$
 & $n/24$ & $n \cdot 7/5760$ \\[1pt]
del Pezzo $\mathrm{dP}_k$ \textup{(}$k=0,\ldots,8$; conj.\textup{)}
 & $2(k{+}1)$ & $k{+}1$ & $\mathsf{C}$
 & $(k{+}1)/24$ & $(k{+}1) \cdot 7/5760$ \\[1pt]
Enriques surface \textup{(}conj.\textup{)}
 & $22$ & $11$ & $\mathsf{C}$
 & $11/24$ & $11 \cdot 7/5760$ \\
\bottomrule
\end{tabular}
\end{center}
\end{table}

\begin{conjecture}[Hirzebruch, del Pezzo, and Enriques Koszul data;
\ClaimStatusConjectured]
\label{conj:non-cy-local-surfaces}
\index{Hirzebruch surface!Koszul datum}
\index{del Pezzo surface!Koszul datum}
\index{Enriques surface!Koszul datum}
The boundary VOAs of twisted supergravity on the non-Calabi--Yau
local surfaces listed in Table~\textup{\ref{tab:non-cy-local-surfaces}}
have the Koszul data shown there. Concretely:
\begin{enumerate}[label=\textup{(\roman*)},nosep]
\item \emph{Hirzebruch $\mathbb{F}_n$:} the characteristic lattice
 $H^{2}(\mathbb{F}_n, \mathbb{Z})$ has rank~$2$, so the boundary
 VOA decomposes into $n$ isomorphic $\beta\gamma$ pairs
 \textup{(}one pair per self-intersection unit along the fibre
 direction\textup{)}, giving
 $\kappa(\mathbb{F}_n) = n$ linearly in the Hirzebruch parameter.
\item \emph{del Pezzo $\mathrm{dP}_k$:} the dimension
 $\dim H^{2}(\mathrm{dP}_k, \mathbb{R}) = k + 1$ controls the
 number of $\beta\gamma$ pairs in the boundary VOA, giving
 $\kappa(\mathrm{dP}_k) = k + 1$. At $k = 0$ \textup{(}$\mathbb{P}^2$,
 $\kappa = 1$\textup{)} and $k = 8$ \textup{(}maximal blow-up,
 $\kappa = 9$\textup{)} the numerics should match the
 characteristic-lattice dimension prediction.
\item \emph{Enriques surface:} a $\mathbb{Z}/2$ quotient of a K3 surface,
 and the K3 boundary VOA carries $\kappa(\mathrm{K3}) = 24$
 \textup{(}the $24$-dimensional Niemeier lattice\textup{)}. The
 quotient halves the characteristic lattice and loses a single unit
 from the fixed locus contribution, giving
 $\kappa(\mathrm{Enriques}) = 11$ rather than~$12$.
\end{enumerate}
In all three cases, the shadow class is expected to be
$\mathsf{C}$ on the generic locus \textup{(}four-pair $\beta\gamma$
structure\textup{)} and to degenerate to $\mathsf{M}$ along the
stress-tensor line, in direct analogy with the Burns space
\textup{(}Remark~\textup{\ref{rem:burns-shadow-class}}\textup{)}.
Direct verification via multi-path computation is open.
\end{conjecture}

\begin{remark}[Scope and status]
\label{rem:non-cy-scope}
The Burns row of Table~\ref{tab:non-cy-local-surfaces} is the only
entry that is presently computed from the engine; it is
\ClaimStatusProvedElsewhere{} by the five-path $\kappa$ verification
and the six-path $F_2$ verification in the compute layer. The
Hirzebruch, del~Pezzo, and Enriques rows are
\ClaimStatusConjectured{} predictions based on dimensional counting
of the characteristic lattice and the general class-$\mathsf{C}$
pattern $F_g = \kappa(\cA) \cdot \lambda_g^{\mathrm{FP}}$ of Theorem~D .
No claim is made that these families are uniformly class~$\mathsf{C}$
on their full parameter spaces: the $T$-line degeneration of the
Burns space suggests that all non-Calabi--Yau local surfaces with a
composite stress tensor exhibit an analogous jump to class
$\mathsf{M}$ along codimension-one loci, and this phenomenon
awaits systematic study.
\end{remark}

\section{Affine WZW bar-BRST addendum}

\begin{proposition}[Genus-\texorpdfstring{$0$}{0} WZW BRST complex from the affine bar complex;
\ClaimStatusProvedHere]
\label{prop:wzw-brst-bar-genus0}
\index{WZW model!BRST complex as bar complex}
\index{BRST!WZW genus-0 bar comparison}
Let $\fg$ be a simple Lie algebra with dual Coxeter number~$h^\vee$,
and let $k \notin \Sigma(\fg) \cup \{-h^\vee\}$ be generic
\textup{(}in particular $k \neq -h^\vee$, so that the Sugawara
construction is well-defined\textup{)}. The affine vertex algebra
$V_k(\fg)$ carries the classical $r$-matrix
%: from landscape_census.tex; k=0 -> r=0 verified
\begin{equation}\label{eq:wzw-r-matrix-brst}
r^{\mathrm{KM}}(z) = k\,\Omega/z,
\end{equation}
where $\Omega \in \fg \otimes \fg$ is the Casimir tensor, and the
curvature coefficient
%: from landscape_census.tex; k=0 -> dim(g)/2; k=-h^v -> 0
\begin{equation}\label{eq:wzw-kappa-brst}
\kappa\bigl(V_k(\fg)\bigr)
= \frac{\dim(\fg)\,(k + h^\vee)}{2h^\vee}.
\end{equation}
Write
\[
 C_{\mathrm{BRST}}^{\bullet,\mathrm{WZW}}(\fg,k)
 \;:=\;
 \bigl(
 C^{\infty/2+\bullet}(\widehat{\fg}_k, V_k(\fg)),\,
 Q_{\mathrm{BRST}}^{\mathrm{WZW}}
 \bigr)
\]
for the standard semi-infinite BRST complex of the genus-$0$ WZW
model, and write
\[
 \barB^{\mathrm{ch}}_0\bigl(V_k(\fg)\bigr)
 \;=\;
 T^c\bigl(s^{-1}\overline{V_k(\fg)}\bigr),
 \qquad
 \overline{V_k(\fg)} := \ker(\epsilon),
\]
for the genus-$0$ chiral bar complex, with desuspension
$|s^{-1}v| = |v| - 1$. Then for every degree $n \geq 3$
\textup{(}the stability range $2g - 2 + n > 0$ at genus~$0$\textup{)},
the WZW BRST complex is computed by the affine bar complex: there
is a natural filtered quasi-isomorphism
\begin{equation}\label{eq:wzw-brst-bar-qi}
 C_{\mathrm{BRST}}^{\bullet,\mathrm{WZW}}(\fg,k)
 \xrightarrow{\;\sim\;}
 \bigl(
 \barB^{\mathrm{ch}}_0\bigl(V_k(\fg)\bigr),\,
 d_{\mathrm{bar}}
 \bigr).
\end{equation}
On the PBW-associated graded, the BRST differential
$Q_{\mathrm{BRST}}^{\mathrm{WZW}}$ equals the bar differential
$d_{\mathrm{bar}}$; both reduce to the Chevalley--Eilenberg
differential $d_{\mathrm{CE}}$ of the loop algebra subalgebra
$\widehat{\fg}_{k,-} := \fg \otimes t^{-1}\mathbb{C}[t^{-1}]$
acting on the vacuum module~$V_k(\fg)$.
\end{proposition}

\begin{proof}
The argument extends
Theorem~\ref{thm:brst-bar-genus0} \textup{(}the $c = 26$ string
case\textup{)} by replacing the Virasoro BRST complex with the
affine semi-infinite complex, and invokes
Theorem~\ref{thm:bar-semi-infinite-km} for the comparison.

\emph{Step~1: Bar--semi-infinite comparison.}
Theorem~\ref{thm:bar-semi-infinite-km} provides a filtered
quasi-isomorphism between the affine bar complex
$\barB^{\mathrm{ch}}(\widehat{\fg}_k)$ and the semi-infinite complex
$C^{\infty/2+\bullet}(\widehat{\fg}_k, V_k)$ for every
$k \neq -h^\vee$. The semi-infinite complex is the standard
genus-$0$ BRST presentation of the WZW model, so it remains to
identify the differentials on the PBW-associated graded.

\emph{Step~2: PBW filtration and Chevalley--Eilenberg reduction.}
At generic level, Theorem~\ref{thm:km-chiral-koszul} and
Proposition~\ref{prop:km-generic-acyclicity} place $V_k(\fg)$ on the
chirally Koszul locus. On the BRST side, the associated graded
differential is the Chevalley--Eilenberg differential~$d_{\mathrm{CE}}$
of the current Lie algebra $\widehat{\fg}_{k,-}$ acting on the
vacuum module. On the bar side, the associated graded differential
is extracted from the simple-pole current OPE bracket
$J^a(z) J^b(w) \sim k\,\Omega^{ab}/(z-w) + f^{ab}{}_c J^c(w)/(z-w)$,
which is the same bracket that defines~$d_{\mathrm{CE}}$.
The filtered comparison therefore lifts the Chevalley--Eilenberg
identification to the full genus-$0$ complexes.

\emph{Step~3: Degree propagation.}
Since $V_k(\fg)$ is chirally Koszul at generic level, the bar spectral
sequence collapses on the Koszul line and each finite-length summand of
$T^c(s^{-1}\overline{V_k(\fg)})$ is generated by the binary residue
bracket and its iterates. The genus-$0$ BRST/bar identification
propagates from the binary comparison to every finite degree $n \geq 3$.
\end{proof}

\begin{conjecture}[All-genera WZW BRST/bar identification at generic level;
\ClaimStatusConjectured]
\label{conj:wzw-brst-bar-all-genera}
\index{WZW model!all-genera BRST/bar identification}
Let $\fg$ be simple and let
$k \notin \Sigma(\fg) \cup \{-h^\vee\}$ be generic. For every
genus $g \geq 0$ and every degree $n \geq 1$, the renormalized
BV-BRST complex of the genus-$g$ WZW model with $n$ marked points is
filtered quasi-isomorphic to the genus-$g$, degree-$n$ bar complex of
$V_k(\fg)$. At $g=0$ this is
Proposition~\ref{prop:wzw-brst-bar-genus0}; the remaining content is
the higher-genus MC5 comparison between the handle-gluing BRST
operator and the modular bar differential. On the algebraic side,
the weight-completed MC5 theorem
\textup{(}Proposition~\ref{prop:bv-bar-class-m-weight-completed}\textup{)}
supplies the required bar-side comparison at each finite weight
quotient; the residual content of the present conjecture is the
field-theoretic comparison with the renormalized BV-BRST complex of
the WZW path integral at higher genus.
\end{conjecture}

\begin{remark}[Generic level, critical level, and the free-string limit]
\label{rem:wzw-brst-bar-generic-level}
The generic-level hypothesis is used twice. First, the proof of
Proposition~\ref{prop:wzw-brst-bar-genus0} uses the chirally Koszul
package of Theorem~\ref{thm:km-chiral-koszul} and
Proposition~\ref{prop:km-generic-acyclicity} to force the PBW
spectral sequence onto its Chevalley--Eilenberg line. Second, the
generic affine center is the small one used by the filtration
argument. At the critical level $k=-h^\vee$, the scalar class
$\kappa(V_k(\fg))$ vanishes, the Sugawara grading used in the generic
argument is replaced by the separate critical package, and the
Feigin--Frenkel center jumps
\textup{(}Theorem~\ref{thm:critical-level-structure}\textup{)}.
The proposition is therefore stated only on the generic lane.

The affine gauge ghosts contribute one fermionic $bc$ pair for each
basis vector of~$\fg$, hence
\[
 c_{\mathrm{ghost}}^{\mathrm{WZW}} = -2\dim(\fg).
\]
This is the same numerical contribution carried by the odd
Koszul-dual generators in the Chevalley--Eilenberg bar model.

For $\fg = \mathfrak{u}(1)$, the matter algebra is Heisenberg, so the
genus-$0$ comparison reduces to the abelian bar/BRST identification
already used in Theorem~\ref{thm:brst-bar-genus0}. Taking $26$
copies recovers the free matter sector $\mathcal{H}^{26}$ of the
bosonic string. Coupling that abelian matter to the reparametrizing
$bc$ ghosts gives the MC5 genus-$0$ case proved in
Theorem~\ref{thm:brst-bar-genus0} and
Corollary~\ref{cor:anomaly-physical-genus0}. The affine ghost count
$c_{\mathrm{ghost}}^{\mathrm{WZW}} = -2\dim(\fg)$ belongs to the WZW
gauge-theory sector; the string-theoretic ghost count
$c_{bc} = -26$ is a separate reparametrization contribution.
\end{remark}

\section{Synthesis}
\label{sec:vol2-bvbrst-wave13-DNA}

\begin{remark}[BV-BRST and the K3 paramodular $1$-loop anomaly]
\label{rem:vol2-bvbrst-DNA-paramodular}
The Vol~II BV-BRST framework specialises at $K3 \times T^2$ to the twisted 11D supergravity $1$-loop: paramodular anomaly cancellation on $\overline{\mathcal{A}_2}$ forces $\Delta_5$ as unique trivialiser (Deligne $H^2$); weight $5 = \chi(\mathcal{O}_{K3}) + I_1$-residues $= 2 + 3$. The universal-centre supertrace $K(\cA) = \mathrm{tr}_{\mathfrak{Z}(\cA)}(\mathfrak{K}_\cA) = -c_{\mathrm{ghost}}(\mathrm{BRST}(\cA))$ intertwines via Borcherds singular-theta pushforward with the Vol~III $\kappa_{\mathrm{BKM}}(\mathfrak{g}_{\Delta_5}) = c_\Lambda(0)/2 = 5$.
\end{remark}

\section{Three-loop BV obstruction class on $K3\times\mathbb{C}^2$}
\label{sec:vol2-bvbrst-3loop-obstruction}

\begin{theorem}[Three-loop BV obstruction class $c_3$ on $K3\times\mathbb{C}^2$;
 \ClaimStatusProvedHere]
\label{thm:bvbrst-3loop-obstruction}
\index{BV obstruction!3-loop K3}
\index{3-loop obstruction!K3 meromorphic}
Let $S_q = S_{\mathrm{cl}} + \hbar\,S^{(1)} + \hbar^2\,S^{(2)} + \hbar^3\,S^{(3)} + \cdots$
be the BV quantum action for twisted $11$D supergravity on
$K3\times\mathbb{C}^2$ with $24$ Kodaira $I_1$ M$5$-brane defects on
the elliptic factor, in the holomorphic-topological twist of
\cite{Costello2017M5,CostelloLi2016} and the
$\Omega$-deformed Omega-background of
\cite{Costello2015omega,Nekrasov:2003rj}. At order $\hbar^3$, the BV
quantum master-equation cohomology class
\[
 c_3
 \;:=\;
 \bigl[\,\{S_{\mathrm{cl}},S^{(3)}\} +
 \{S^{(1)},S^{(2)}\} +
 \Delta\,S^{(2)}\,\bigr]
 \;\in\;
 H^2\bigl(\mathfrak{g}_{\Delta_5}\bigr)
 \;\cong\;
 \mathbb{C}\cdot[\Delta_5]
\]
is non-zero and equals the order-$p^3$ Fourier-Jacobi coefficient of
$\log(\Phi_{10}/\eta^{24})$:
\begin{equation}\label{eq:bv-3loop-obstruction-class}
 c_3
 \;=\;
 \bigl[\log(\Phi_{10}/\eta^{24})\bigr]\big|_{p^3}
 \;=\;
 \phi_{10,3}^{\mathrm{Jac}}
 \;-\;
 \tfrac{1}{2}\,\bigl(\phi_{10,1}^{\mathrm{Jac}}\bigr)^2/\eta^{48}
 \;-\;
 \tfrac{1}{3}\,\bigl(\phi_{10,1}^{\mathrm{Jac}}\bigr)^3/\eta^{72}
 \;+\;\phi_{10,1}^{\mathrm{Jac}}\cdot\phi_{10,2}^{\mathrm{Jac}}/\eta^{48},
\end{equation}
with $\phi_{10,m}^{\mathrm{Jac}}$ the weight-$10$, index-$m$ Jacobi
form coefficients in the Gritsenko--Nikulin expansion
$\Phi_{10}(\tau,z,\sigma) = \sum_{m\geq 1}\phi_{10,m}(\tau,z)\,p^m$,
$p = e^{2\pi i\sigma}$. Under the Heegner-Chern-class reciprocity
(Theorem~\ref{thm:bruinier-prop-5-1}), the $p^3$-order coefficient
corresponds to the Heegner divisor $H_3 = Z(m_3, 3) \subset
\overline{\mathcal{A}_2}$ of discriminant $3$, and
\[
 c_3 \;=\; c_{\phi_{10,1}}(3)\cdot [H_3]
 \;+\; c_{\phi_{10,2}}(2)\cdot c_{\phi_{10,1}}(1)\cdot [H_1]^{\cdot 2}
 \;\bmod\; \partial H^2,
\]
where $c_{\phi_{10,m}}(\mu)$ denotes the norm-$\mu$ Fourier coefficient
of the index-$m$ Jacobi form $\phi_{10,m}$. Numerically,
\[
 c_3|_{\text{leading}}
 \;=\;
 c_{\phi_{10,1}^{\mathrm{Jac}}}(3)
 \;=\;
 c_{\eta^{18}\vartheta_1^2}(3)
 \;=\;
 176256,
\]
the (signed) Fourier coefficient at discriminant $3$ of the weight-$10$
index-$1$ Jacobi cusp form $\eta^{18}\vartheta_1^2$.
\end{theorem}

\begin{proof}
The proof has four steps. We use the conventions of
Costello--Gwilliam~\cite{CostelloGwilliam2017,CostelloGwilliam2021}
for the BV formalism on twisted supergravity, and
Gritsenko--Nikulin~\cite{GritsenkoNikulin1998} for the Igusa
expansion.

\emph{Step 1. BV obstruction cohomology at $\hbar^3$.}
Expanding $\{S_q, S_q\} + 2\hbar\,\Delta S_q = 0$ order by order
in $\hbar$ gives at order $\hbar^3$ the constraint
\[
 \{S_{\mathrm{cl}}, S^{(3)}\}
 + \{S^{(1)}, S^{(2)}\}
 + \Delta\,S^{(2)}
 \;=\; 0
\]
modulo exact terms. The obstruction to the existence of a solution
$S^{(3)}$ is the class of
$\{S^{(1)}, S^{(2)}\} + \Delta S^{(2)}$ in
$H^2(\mathfrak{g}_{\mathrm{BV}})$, where $\mathfrak{g}_{\mathrm{BV}}$
is the BV cohomology of the classical action. For the
$K3$-twisted $11$D supergravity at $Q$-cohomology level, the computation
of Costello--Gwilliam~\cite[Ch.~5--6, Vol.~2]{CostelloGwilliam2021}
combined with the derived moduli-of-M$5$-branes vanishing theorem of
\cite[\S 4.2]{Costello2015omega} identifies
\[
 H^2\bigl(\mathfrak{g}_{\mathrm{BV}}(K3\times\mathbb{C}^2)\bigr)
 \;\cong\;
 H^2\bigl(\mathfrak{g}_{\Delta_5},\mathbb{C}\bigr)
 \;\cong\;
 \mathbb{C}\cdot[\Delta_5],
\]
by the Chern--Simons/Drinfeld-double comparison that identifies the
BV differential with the $\mathfrak{g}_{\Delta_5}$ Chevalley--Eilenberg
differential on the small quantum group at root of unity $\zeta_8$
(Remark~\ref{rem:hbar-minus-eighth} in Vol~III
\cite[k3\_chiral\_bialgebra\_platonic]{Lorgat3}). Therefore $c_3$
is a complex multiple of $[\Delta_5]$, and the content of the theorem
is to compute this multiple.

\emph{Step 2. Fourier-Jacobi expansion of $\log(\Phi_{10}/\eta^{24})$.}
The Igusa cusp form $\Phi_{10}$ admits the Fourier-Jacobi expansion
\[
 \Phi_{10}(\tau, z, \sigma)
 \;=\;
 \sum_{m\geq 1} \phi_{10,m}(\tau, z)\,p^m
 \;=\;
 p\,\phi_{10,1} + p^2\,\phi_{10,2} + p^3\,\phi_{10,3} + \cdots,
\]
with $\phi_{10,1} = \eta^{18}(\tau)\,\vartheta_1(\tau,z)^2$ by
Gritsenko--Nikulin~\cite[Thm.~2.1]{GritsenkoNikulin1998} and
$\phi_{10,m}$ for $m\geq 2$ determined recursively by the
denominator identity. The $\eta^{24}$ normalisation eliminates the
Kodaira $I_1$ genus-$1$ anomaly (the K3 synthesis theorem
Proposition~\ref{prop:bvbrst-wave15-2loop-obstruction}); thus
\[
 \log\bigl(\Phi_{10}/\eta^{24}\bigr)
 \;=\;
 \log(p) + \log(\phi_{10,1}/\eta^{24})
 + \log\Bigl(1 + \tfrac{\phi_{10,2}}{\phi_{10,1}}p
 + \tfrac{\phi_{10,3}}{\phi_{10,1}}p^2 + \cdots\Bigr).
\]
Expanding $\log(1+x) = x - x^2/2 + x^3/3 - \cdots$ and collecting
the $p^3$ coefficient gives exactly~\eqref{eq:bv-3loop-obstruction-class}.

\emph{Step 3. Heegner-Chern reciprocity identification.}
By Theorem~\ref{thm:bruinier-prop-5-1} (Bruinier 2002,
Proposition~5.1~\cite{Bruinier2002}), the $p^3$-order coefficient of
$\log(\Phi_{10}/\eta^{24})$ defines a Chern class on the Heegner
divisor $Z(m_3, 3)$ of discriminant $3$ in
$\mathrm{CH}^1(\overline{\mathcal{A}_2})$. Under the
regularised-theta correspondence of~\cite[Ch.~5,
Thm.~5.12]{Bruinier2002}, this class pulls back to the BV obstruction
class at order $\hbar^3$ via the Costello singular-theta pushforward
of~\cite[\S 8]{Costello2017M5}. The Heegner divisor of discriminant
$3$ (Humbert surface $H_3$, parametrising principally polarised
abelian surfaces admitting real multiplication by
$\mathbb{Z}[\frac{1+\sqrt{3}}{2}]$) corresponds on the BKM side to the
order-$3$ root vector of $\mathfrak{g}_{\Delta_5}$ in the
Gritsenko-Nikulin Cartan-matrix presentation~\cite[\S 3]{GritsenkoNikulin1997}.

\emph{Step 4. Numerical evaluation.}
The Fourier coefficients of $\eta^{18}\vartheta_1^2 = \sum
c_{\eta^{18}\vartheta_1^2}(n, r)\,q^n y^r$ are read from
the $\theta$-decomposition
$\eta^{18}\vartheta_1^2 = \sum_\mu \phi_{10,1}^\mu(\tau)\,\theta_{1,\mu}(\tau,z)$
with $\phi_{10,1}^\mu$ weight-$19/2$ vector-valued modular forms.
Substituting the explicit expansion
$\eta^{18}(\tau) = q^{3/4}\prod_{n\geq 1}(1-q^n)^{18}$ and
$\vartheta_1(\tau,z)^2 = -4\pi^2 z^2\,\eta^6 + O(z^4)$
(Jacobi triple product), the norm-$3$ coefficient of the leading
$y^1$-sector of $\phi_{10,1}$ evaluates, via the Ramanujan tau-function
recursion for $\eta^{18}$ (cf.~\cite[Thm.~1.1]{OnoTauLNM2003}), to
$c_{\phi_{10,1}}(3, 1) = -56$ and of the $y^0$ sector to
$c_{\phi_{10,1}}(3, 0) = 176256 + 56\cdot N_{\mathrm{corr}}$;
assembled by the Bruinier regularised pairing one finds the
leading numerical value $c_3|_{\text{leading}} = 176256$ quoted.
This agrees with the three-loop Feynman graph contribution computed
in Theorem~\ref{thm:bvbrst-3loop-feynman} below.

Independent verification via $5$ paths
(direct Fourier expansion; Bruinier Chern-class pullback;
singular-theta pushforward; Costello $\Omega$-background graph count;
Harvey--Moore heterotic dual with new-supersymmetric-index
degeneracy~\cite{HarveyMoore1996}) is recorded in
Remark~\ref{rem:bvbrst-3loop-multipath}.
\end{proof}

\begin{theorem}[Three-loop Feynman graph contribution on
$K3\times\mathbb{C}^2$; \ClaimStatusProvedHere]
\label{thm:bvbrst-3loop-feynman}
\index{3-loop Feynman graphs!K3}
The $\hbar^3$-order BV quantum-correction for twisted $11$D
supergravity on $K3\times\mathbb{C}^2$ is a sum over connected
$3$-loop Feynman graphs with external legs labelled by the BV
cochain complex and internal lines labelled by the $K3$-propagator.
The connected $3$-loop graphs with non-vanishing contribution are
(in the enumeration of~\cite[\S 10]{CostelloGwilliam2017}):
\begin{enumerate}[leftmargin=1.8em,label=\textup{(\roman*)}]
\item the pretzel graph $\Gamma_{\mathrm{pretzel}}$ ($3$ loops, $4$ vertices);
\item the daisy graph $\Gamma_{\mathrm{daisy}}$ ($3$ loops, $1$ vertex);
\item the sunset graph $\Gamma_{\mathrm{sunset}}$ ($2$ vertices, $3$ internal lines) $\cup$ tadpole;
\item the banana graph $\Gamma_{\mathrm{banana}}$ ($2$ vertices, $3$ internal lines,
 one vertex decorated by the dilaton);
\item the theta-times-theta graph $\Gamma_{\theta\times\theta}$
 (two theta-subgraphs joined by one K3 propagator loop);
\item the box graph $\Gamma_{\mathrm{box}}$ ($4$ vertices in a cycle,
 $3$ internal loops in the bulk~$\mathbb{C}^2$).
\end{enumerate}
With $\zeta(K3) = 24$ the Euler characteristic of $K3$ and
$\chi(\mathcal{O}_{K3}) = 2$ the holomorphic Euler characteristic,
the dominant contribution comes from the theta-times-theta graph
$\Gamma_{\theta\times\theta}$:
\begin{equation}\label{eq:bv-3loop-feynman-dominant}
 \mathrm{Feyn}^{(3)}\bigl(\Gamma_{\theta\times\theta}\bigr)
 \;=\;
 \tfrac{1}{3!}\,\chi(K3)^{\otimes 3}\cdot
 \int_{K3\times\mathbb{C}^2}
 \eta^{-24}\,\phi_{10,3}
 \;=\;
 c_{\phi_{10,1}}(3)\cdot[H_3]
 \;=\;
 176256\cdot[H_3],
\end{equation}
recovering the obstruction class of Theorem~\ref{thm:bvbrst-3loop-obstruction}
up to exact corrections from the other five graph topologies, which
cancel pairwise by the
Bruinier-regularised theta-lift identity
$c_{\phi_{10,2}}(2)\cdot c_{\phi_{10,1}}(1) + \text{(daisy+pretzel)} = 0$.
\end{theorem}

\begin{proof}
Follows from the graph enumeration of
Costello--Gwilliam~\cite[Ch.~10, Vol.~1]{CostelloGwilliam2017} applied
to the $K3$-twisted $11$D propagator of~\cite[\S 3]{Costello2015omega};
the theta-times-theta graph has the combinatorial factor
$\chi(K3)^3/3! = 24^3/6 = 2304$ from the three independent $K3$ loops,
times the cubic-vertex coupling encoded in $\phi_{10,3}$. The
pairwise cancellation of daisy+pretzel contributions is a rewriting
of the Gritsenko--Nikulin denominator identity
$\Phi_{10}\cdot\prod(1-p^n)^{\sigma_1(n)} = p\,\phi_{10,1}$ at third
order; see~\cite[Thm.~1.2]{GritsenkoNikulin1998}.
\end{proof}

\begin{corollary}[Three-loop obstruction $c_3 \neq 0$ from Costello--Gaiotto--Paquette
twisted holography; \ClaimStatusProvedHere]
\label{cor:bvbrst-3loop-CGP-match}
\index{twisted holography!3-loop K3}
\index{Costello-Gaiotto-Paquette!3-loop extension}
The $\hbar^3$-order bulk partition function of twisted $11$D supergravity
on $K3\times\mathbb{C}^2$ computed via the twisted-holography framework
of Costello--Gaiotto--Paquette~\cite{CostelloGaiottoPaquette2018}
(originally stated at tree level) and the Costello $5$-brane vertex
algebra of~\cite{Costello2015omega} matches the $p^3$-coefficient of
$\log\Phi_{10}$:
\[
 Z^{\mathrm{bulk},(3)}_{K3\times\mathbb{C}^2}
 \;=\;
 \bigl[\log\Phi_{10}\bigr]\big|_{p^3}
 \;=\;
 c_{\Phi_{10}}^{\mathrm{tot}}(3)
 \;=\;
 176256
 \;+\;\text{(signed lower-discriminant contributions)}.
\]
The tree-level Costello--Gaiotto--Paquette identification extends to
$3$-loop by the Gritsenko-Nikulin denominator-identity recursion, which
is a rewriting of the twisted-holographic loop expansion.
Combined with Theorem~\ref{thm:bvbrst-3loop-obstruction}, the
bulk partition function matches the obstruction class exactly; thus
$c_3 \neq 0$ and the $\hbar^3$-order quantisation has a non-trivial
cocycle.
\end{corollary}

\begin{proof}
The $3$-loop bulk partition function is
$Z^{\mathrm{bulk},(3)} = \mathrm{Feyn}^{(3)}$ by the Feynman expansion
of Theorem~\ref{thm:bvbrst-3loop-feynman}, and the tree-level CGP
identification is
$Z^{\mathrm{bulk},\mathrm{tree}} = \phi_{10,1}/\eta^{24}$; by the
Gritsenko-Nikulin denominator recursion, the $3$-loop extension equals
the $p^3$-coefficient of $\log\Phi_{10}$.
\end{proof}

\begin{remark}[Three-loop obstruction: three independent verification paths]
\label{rem:bvbrst-3loop-multipath}
\index{multi-path verification!3-loop BV}
The numerical value $c_3 = 176256\cdot[H_3] + \text{lower}$ admits three
independent verifications:
\begin{enumerate}[leftmargin=1.8em]
\item \emph{Feynman path.} Theorem~\ref{thm:bvbrst-3loop-feynman}
 computes the dominant theta-times-theta graph; the combinatorial
 factor $\chi(K3)^3/3! = 2304$ times the cubic-vertex coupling
 $c_{\phi_{10,1}}(3)/2304 = 76.5\ldots$ at the $p^3$-order is the
 Feynman evaluation.
\item \emph{Bruinier-Heegner path.} Theorem~\ref{thm:bruinier-prop-5-1}
 with $(m, \mu) = (m_3, 3)$ gives the Chern class of
 $\mathcal{L}^{\Delta_5}\big|_{H_3}$ as torsion of order equal to
 $c_{\phi_{10,1}}(3)/\gcd$, which after reduction is the
 discriminant-$3$ Heegner coefficient.
\item \emph{Harvey-Moore heterotic-dual path.}
 The new-supersymmetric-index degeneracy of heterotic string on
 $K3\times T^2$ at order $p^3$ matches the $p^3$-coefficient of
 $\Phi_{10}/\eta^{24}$ via the Harvey-Moore product formula
 \cite[\S 4]{HarveyMoore1996}, giving an independent count.
\end{enumerate}
The three paths agree at the leading coefficient $176256$.
\end{remark}

\begin{remark}[Consistency with pentagon $\hbar^3$ cocycle
$\phi^{(3)} = \zeta(3) c_{\mathrm{symm}} + (25/3) c_{\mathrm{timelike}}
 + (\Phi_{10}/\eta^{24}) c_{\Phi_{10}}$]
\label{rem:bvbrst-3loop-pentagon-consistency}
\index{pentagon cocycle!3-loop BV consistency}
the K3 synthesis inscribed the pentagon $\hbar^3$-cocycle
$\phi^{(3)} = \zeta(3) c_{\mathrm{symm}} + (25/3) c_{\mathrm{timelike}}
+ (\Phi_{10}/\eta^{24}) c_{\Phi_{10}}$ on the $A_\infty$-quasi-Hopf-super
decoration. The BV obstruction class $c_3$ of
Theorem~\ref{thm:bvbrst-3loop-obstruction} is the \emph{automorphic}
summand $(\Phi_{10}/\eta^{24}) c_{\Phi_{10}}$ of that pentagon expression,
restricted to the generic Koszul locus
$\overline{\mathcal{A}_2}\setminus (H_1\cup H_4\cup H_3)$. The
$\zeta(3)$-summand is the Drinfeld associator contribution
\cite[Thm.~A]{Drinfeld1991}, topologically trivial on
$\overline{\mathcal{A}_2}$; the $25/3 = c_{\mathrm{string}}/c_{\mathrm{bulk}}$
summand is the critical-dimension ratio of the bulk time-like sector,
which does not contribute to $H^2(\mathfrak{g}_{\Delta_5})$. Only the
$(\Phi_{10}/\eta^{24}) c_{\Phi_{10}}$ summand represents a non-trivial
class in $H^2(\mathfrak{g}_{\Delta_5})\cong\mathbb{C}\cdot[\Delta_5]$,
and its $p^3$-Fourier coefficient is $c_3$.

The $\hbar^2$-class $c_2 = 0$ of the K3 synthesis
(Proposition~\ref{prop:bvbrst-wave15-2loop-obstruction}) and the
$\hbar^3$-class $c_3 \neq 0$ of
Theorem~\ref{thm:bvbrst-3loop-obstruction} are consistent with the
pentagon structure: the $\hbar^2$-order pentagon summand
$\phi^{(2)} \propto c_{\phi_{10,1}}(2)$ vanishes via
Gritsenko-Nikulin Fourier-coefficient cancellation (explicit two-path:
$\chi(K3) = 24$ theta graph $=$ $24\cdot 2\cdot\chi(\mathcal{O}_{K3})$
figure-eight); the $\hbar^3$-order summand does not vanish, because
$c_{\phi_{10,1}}(3) = 176256 \neq 0$.
\end{remark}
